%% toc: P3 31 31.1,31.2,31.3,31.4
\section{Proofs for indiscernibles and Ehrenfeucht--Mostowski models}
\label{sec:annex-p03-w03-p03-028}

The proofs of Theorems~A, B and C follow Ramsey's arguments
\cite[pp.~264--269]{Ramsey1930}, with the case distinctions written out.
Throughout, ``class'' and ``set'' are used interchangeably.

\begin{annexitem}{annex:p03-w03-p03-028-ramsey-a}{thm:p03-w03-p03-028-ramsey-a}{Ramsey's Theorem A}
\begin{proof}
If \(\mu=1\) there is nothing to prove.  First let \(\mu=2\), with classes
\(C_1\) and \(C_2\).  We prove the claim for every \(r\geq1\) by induction
on \(r\).

\emph{Case \(r=1\).}  Each member \(x\) of \(\Gamma\) lies in
\(\{y:\{y\}\in C_1\}\) or in \(\{y:\{y\}\in C_2\}\).  These two sets cover
the infinite set \(\Gamma\), so one of them is infinite; it is the required
\(\Delta\).

\emph{Inductive step.}  Let \(\rho\geq2\) and assume the claim for
\(r=\rho-1\) and \(\mu=2\), for every infinite class.  Call a pair
\((x,\Gamma')\) \emph{good for \(C_1\) inside} a class \(\Lambda\) when
\(x\in\Lambda\), \(\Gamma'\) is an infinite subclass of
\(\Lambda\setminus\{x\}\), and every \(\rho\)-combination consisting of
\(x\) and \(\rho-1\) members of \(\Gamma'\) lies in \(C_1\).  Put
\(\Gamma_0=\Gamma\) and try to choose, for \(n=1,2,\ldots\), a pair
\((x_n,\Gamma_n)\) good for \(C_1\) inside \(\Gamma_{n-1}\).

\emph{First possibility: the selection never stops.}  Then
\(\Gamma_0\supseteq\Gamma_1\supseteq\cdots\) and \(x_n\in\Gamma_{n-1}\),
\(x_n\notin\Gamma_n\).  The \(x_n\) are distinct: if \(t>s\) then
\(x_t\in\Gamma_{t-1}\subseteq\Gamma_s\), while \(x_s\notin\Gamma_s\).  Let
\(\Delta=\{x_1,x_2,\ldots\}\), an infinite class.  Let \(R\) be a
\(\rho\)-combination of \(\Delta\) and let \(x_s\) be its member of least
suffix.  Its other \(\rho-1\) members have larger suffixes, so they lie in
\(\Gamma_s\).  Since \((x_s,\Gamma_s)\) is good for \(C_1\), \(R\in C_1\).

\emph{Second possibility: the selection stops.}  Then for some \(n\) the
class \(\Lambda=\Gamma_{n-1}\) is infinite and has no pair good for \(C_1\)
inside it.  We choose \(y_1,y_2,\ldots\) and infinite classes
\(\Lambda\supseteq\Delta_1\supseteq\Delta_2\supseteq\cdots\) as follows.
Let \(\Delta_0=\Lambda\).  Given an infinite \(\Delta_{t-1}\subseteq
\Lambda\), pick \(y_t\in\Delta_{t-1}\).  Divide the
\((\rho-1)\)-combinations \(Q\) of \(\Delta_{t-1}\setminus\{y_t\}\) into
two classes according as \(Q\cup\{y_t\}\) lies in \(C_1\) or in \(C_2\).
By the induction hypothesis there is an infinite
\(\Delta_t\subseteq\Delta_{t-1}\setminus\{y_t\}\) all of whose
\((\rho-1)\)-combinations fall in one of these classes, so that all
\(\rho\)-combinations consisting of \(y_t\) and \(\rho-1\) members of
\(\Delta_t\) lie in one \(C_i\).  That class is not \(C_1\): otherwise
\((y_t,\Delta_t)\) would be good for \(C_1\) inside \(\Lambda\), because
\(y_t\in\Lambda\) and \(\Delta_t\subseteq\Lambda\setminus\{y_t\}\).  So
it is \(C_2\).  As in the first possibility, the \(y_t\) are distinct, and
for \(\Delta=\{y_1,y_2,\ldots\}\) every \(\rho\)-combination, split at its
member \(y_s\) of least suffix, consists of \(y_s\) and \(\rho-1\) members
of \(\Delta_s\), so it lies in \(C_2\).

In both cases \(\Gamma\) has an infinite subclass all of whose
\(\rho\)-combinations lie in one class.  This proves the case \(\mu=2\)
for all \(r\).  The infinitely many successive choices of elements and
infinite subclasses use the axiom of choice, which the statement assumes.

\emph{More than two classes.}  Let \(\nu>2\) and assume the theorem for
\(\mu=2\) and \(\mu=\nu-1\).  Given classes \(C_1,\ldots,C_\nu\), put
\(C'_i=C_i\) for \(i\leq\nu-2\) and \(C'_{\nu-1}=C_{\nu-1}\cup C_\nu\).
These \(\nu-1\) classes divide the \(r\)-combinations, so some infinite
\(\Delta\subseteq\Gamma\) has all its \(r\)-combinations in one \(C'_i\).
If \(i\leq\nu-2\) we are done.  Otherwise all of them lie in
\(C_{\nu-1}\cup C_\nu\); dividing the \(r\)-combinations of \(\Delta\) into
those in \(C_{\nu-1}\) and those in \(C_\nu\), the case \(\mu=2\) gives an
infinite \(\Delta'\subseteq\Delta\) all of whose \(r\)-combinations lie in
\(C_{\nu-1}\) or all in \(C_\nu\).
\end{proof}
\end{annexitem}

\begin{annexitem}{annex:p03-w03-p03-028-bounds-defined}{thm:p03-w03-p03-028-bounds-defined}{The bound functions are well defined}
\begin{proof}
We use induction on \(r\).  For \(r=1\), \(f(1,n,k)=\max(2n-1,n+k)\) is
defined for \(n\geq1\), \(k\geq0\) and is at least \(n+k\).

Let \(r\geq2\) and assume that \(f(r-1,n',k')\) is defined, and at least
\(n'+k'\), whenever \(n'\geq1\), \(k'\geq0\), and either \(n'=1\) and
\(k'+1\geq r-1\), or \(n'>1\) and \(n'+k'\geq r-1\) (for \(r-1=1\) this
covers all \(n'\geq1\), \(k'\geq0\)).

\emph{The clause for \(f(r,1,k)\).}  Let \(k+1\geq r\).  The arguments
\(n'=k-r+2\) and \(k'=r-2\) satisfy \(n'\geq1\) and \(k'\geq0\).  If
\(n'=1\) then \(k'+1=r-1\); if \(n'>1\) then \(n'+k'=k\geq r-1\).  So
\(f(r-1,n',k')\) is defined, and
\(f(r,1,k)\geq(k-r+2)+(r-2)+1=k+1\).

\emph{The function \(g\).}  \(g(r,0,k)=\max(r-1,k)\).  If
\(g(r,n-1,k)\) is defined and at least \(\max(r-1,k)+n-1\), then
\(g(r,n-1,k)+1\geq r\), so \(g(r,n,k)=f(r,1,g(r,n-1,k))\) is defined by the
previous clause and is at least \(g(r,n-1,k)+1\geq\max(r-1,k)+n\).  By
induction on \(n\), \(g(r,n,k)\) is defined for all \(n,k\geq0\).  Only
values \(f(r,1,\cdot)\) are used, so this does not presuppose \(f(r,n,k)\)
for \(n>1\).

\emph{The function \(f(r,n,k)\) for \(n>1\).}  By induction on \(n\geq2\),
assume \(f(r,n-1,k')\) defined and at least \(n-1+k'\) for every \(k'\)
with \(k'+1\geq r\) (this is the case \(n-1=1\)) or with
\(n-1+k'\geq r\) (the case \(n-1>1\)).  The argument
\(k'=g(r,n,k)\geq r-1\) meets both conditions, so
\(f(r,n,k)=f(r,n-1,g(r,n,k))\) is defined for every \(k\geq0\) and is at
least \(n-1+g(r,n,k)\geq n-1+k+n\geq n+k\).

\emph{The function \(h\).}  For \(n\geq r\), \(h(r,n,2)=f(r,n-r+1,r-1)\):
here \(n-r+1\geq1\), and if \(n-r+1=1\) then \((r-1)+1=r\), while if
\(n-r+1>1\) then \((n-r+1)+(r-1)=n\geq r\); so the value is defined and is
at least \(n\).  If \(h(r,n,\mu-1)\) is defined and at least \(n\geq r\),
then \(h(r,n,\mu)=h(r,h(r,n,\mu-1),2)\) is defined and at least
\(h(r,n,\mu-1)\geq n\).  For \(n<r\) or \(\mu=1\), \(h(r,n,\mu)=n\).
\end{proof}
\end{annexitem}

\begin{annexitem}{annex:p03-w03-p03-028-ramsey-c}{thm:p03-w03-p03-028-ramsey-c}{Ramsey's Theorem C}
\begin{proof}
By \cref{thm:p03-w03-p03-028-bounds-defined}, \(m_0=f(r,n,k)\) is defined.
We prove the theorem by induction on \(r\), and for fixed \(r\geq2\) by
induction on \(n\).  Throughout, \(\Gamma_m\) has \(m\geq m_0\) members and
its \(r\)-combinations are divided into \(C_1\) and \(C_2\).

\emph{Case \(r=1\).}  Here \(m\geq\max(2n-1,n+k)\).  If both
\(\{x:\{x\}\in C_1\}\) and \(\{x:\{x\}\in C_2\}\) had fewer than \(n\)
members, \(\Gamma_m\) would have at most \(2n-2\) members.  So one of them
contains some \(\Delta_n\).  Since \(m-n\geq k\), a \(\Delta_k\) can be
chosen in \(\Gamma_m\setminus\Delta_n\).  The \(1\)-combinations of
\(\Delta_n+\Delta_k\) that meet \(\Delta_n\) are the singletons of
\(\Delta_n\), all in one class.

Now let \(r\geq2\) and assume the theorem for \(r-1\) and all admissible
\(n,k\).

\emph{Case \(n=1\).}  Here \(k+1\geq r\) and
\(m\geq f(r-1,k-r+2,r-2)+1\).  Choose any \(x\in\Gamma_m\).  Then
\(\Gamma_m\setminus\{x\}\) has at least \(f(r-1,k-r+2,r-2)\) members.
Divide its \((r-1)\)-combinations \(Q\) into \(C'_1\) and \(C'_2\)
according as \(Q\cup\{x\}\) lies in \(C_1\) or \(C_2\).  The theorem for
\(r-1\), with \(k-r+2\) and \(r-2\) in place of \(n\) and \(k\) (admissible
by \cref{thm:p03-w03-p03-028-bounds-defined}), gives disjoint
\(\Delta_{k-r+2}\) and \(\Delta_{r-2}\) in \(\Gamma_m\setminus\{x\}\) such
that every \((r-1)\)-combination of their union that meets
\(\Delta_{k-r+2}\) lies in one \(C'_i\).  Since \(\Delta_{r-2}\) has only
\(r-2\) members, every \((r-1)\)-combination of the union meets
\(\Delta_{k-r+2}\).  Let \(\Delta_k\) be the union, which has \(k\)
members, and \(\Delta_1=\{x\}\).  An \(r\)-combination of
\(\Delta_1+\Delta_k\) meeting \(\Delta_1\) is \(\{x\}\cup Q\) with \(Q\) an
\((r-1)\)-combination of \(\Delta_k\), so all of them lie in the same
\(C_i\).

\emph{Case \(n>1\).}  Assume the theorem for \(r\) and \(n-1\) with every
admissible \(k\).  Put \(G=g(r,n,k)\), so \(m\geq f(r,n-1,G)\).  The theorem
for \(n-1\) and \(G\) gives disjoint \(\Delta_{n-1}\) and \(\Delta_G\) such
that every \(r\)-combination of \(\Delta_{n-1}+\Delta_G\) that meets
\(\Delta_{n-1}\) lies in one class; call it \(C_a\), and call the other
class \(C_b\).

\emph{Subcase (a).}  Suppose \(\Delta_G\) contains a member \(x\) and a
subclass \(\Delta_k\) not containing \(x\) such that every combination of
\(x\) with \(r-1\) members of \(\Delta_k\) lies in \(C_a\).  Let
\(\Delta_n=\Delta_{n-1}\cup\{x\}\).  An \(r\)-combination \(R\) of
\(\Delta_n+\Delta_k\) meeting \(\Delta_n\) either meets \(\Delta_{n-1}\),
and then lies in \(\Delta_{n-1}+\Delta_G\) and so in \(C_a\); or it consists
of \(x\) and \(r-1\) members of \(\Delta_k\), and then it lies in \(C_a\)
by the choice of \(x\).

\emph{Subcase (b).}  Otherwise no \(x\in\Delta_G\) has such a \(\Delta_k\)
inside \(\Delta_G\setminus\{x\}\).  Put \(\Lambda_0=\Delta_G\), which has
\(g(r,n,k)\) members.  For \(j=1,\ldots,n\) suppose \(\Lambda_{j-1}
\subseteq\Delta_G\) has \(g(r,n-j+1,k)\) members.  Because
\(g(r,n-j+1,k)=f(r,1,g(r,n-j,k))\) and \(g(r,n-j,k)+1\geq r\), the case
\(n=1\) applied to \(\Lambda_{j-1}\) gives \(x_j\in\Lambda_{j-1}\) and
\(\Lambda_j\subseteq\Lambda_{j-1}\setminus\{x_j\}\) with
\(g(r,n-j,k)\) members such that all combinations of \(x_j\) with
\(r-1\) members of \(\Lambda_j\) lie in one class.  That class is not
\(C_a\): since \(g(r,n-j,k)\geq k\), any \(k\) members of \(\Lambda_j\)
would together with \(x_j\) satisfy subcase (a), as \(x_j\in\Delta_G\) and
\(\Lambda_j\subseteq\Delta_G\setminus\{x_j\}\).  So it is \(C_b\).

The \(x_j\) are distinct, since \(x_l\in\Lambda_{l-1}\subseteq\Lambda_j\)
for \(l>j\) and \(x_j\notin\Lambda_j\).  Let
\(\Delta_n=\{x_1,\ldots,x_n\}\) and let \(\Delta_k\) be any \(k\) members of
\(\Lambda_n\), which has \(g(r,0,k)\geq k\) members.  Let \(R\) be an
\(r\)-combination of \(\Delta_n+\Delta_k\) meeting \(\Delta_n\), and let
\(x_j\) be its member of least suffix.  The other members of \(R\) lie in
\(\{x_{j+1},\ldots,x_n\}\cup\Lambda_n\subseteq\Lambda_j\).  So \(R\) consists
of \(x_j\) and \(r-1\) members of \(\Lambda_j\), and \(R\in C_b\).

In both subcases \(\Gamma_m\) contains the required \(\Delta_n\) and
\(\Delta_k\).  This completes both inductions.
\end{proof}
\end{annexitem}

\begin{annexitem}{annex:p03-w03-p03-028-ramsey-b}{thm:p03-w03-p03-028-ramsey-b}{Ramsey's Theorem B}
\begin{proof}
Let \(m\geq m_0=h(r,n,\mu)\).  If \(\mu=1\), or if \(n<r\), any \(n\)
members of \(\Gamma_m\) form a homogeneous \(\Delta_n\): in the first case
there is only one class, in the second there are no \(r\)-combinations of
\(\Delta_n\).  Assume \(n\geq r\).

\emph{Case \(\mu=2\).}  Here \(m\geq f(r,n-r+1,r-1)\), and
\cref{thm:p03-w03-p03-028-ramsey-c} applies with \(n-r+1\) and \(r-1\) in
place of \(n\) and \(k\) (admissible by
\cref{thm:p03-w03-p03-028-bounds-defined}).  It gives disjoint
\(\Delta_{n-r+1}\) and \(\Delta_{r-1}\) such that all \(r\)-combinations of
their union that meet \(\Delta_{n-r+1}\) lie in one class.  Every
\(r\)-combination of the union meets \(\Delta_{n-r+1}\), because
\(\Delta_{r-1}\) has only \(r-1\) members.  So the union, which has \(n\)
members, is the required \(\Delta_n\).

\emph{Case \(\mu>2\).}  Assume the theorem for \(\mu-1\).  Put
\(C'_1=C_1\) and \(C'_2=C_2\cup\cdots\cup C_\mu\).  Since
\(m\geq h(r,h(r,n,\mu-1),2)\), the case \(\mu=2\) gives a subclass
\(\Lambda\) with \(h(r,n,\mu-1)\) members all of whose
\(r\)-combinations lie in \(C'_1\), or all in \(C'_2\).  In the first case
any \(n\) members of \(\Lambda\) form a homogeneous \(\Delta_n\), since
\(h(r,n,\mu-1)\geq n\) by \cref{thm:p03-w03-p03-028-bounds-defined}.  In
the second case the \(r\)-combinations of \(\Lambda\) are divided among the
\(\mu-1\) classes \(C_2,\ldots,C_\mu\), and the theorem for \(\mu-1\)
gives \(\Delta_n\subseteq\Lambda\) homogeneous for one of them.
\end{proof}
\end{annexitem}

\begin{annexitem}{annex:p03-w03-p03-028-ramsey-step}{thm:p03-w03-p03-028-ramsey-step}{The partition step}
\begin{proof}
If \(q=0\) take \(Y=\dom\mathcal M\).  Otherwise let
\(\Gamma=\dom\mathcal M\), \(r=N\) and \(\mu=2^q\), and index the classes
by \(\varepsilon\in\{0,1\}^q\).  For an \(N\)-combination \(U\) of
\(\Gamma\) let \(\bar u\) be its \(\lhd\)-increasing enumeration and put
\(U\in C_\varepsilon\) where \(\varepsilon_j=1\) iff
\(\mathcal M\entails\psi_j(\bar u)\).  Each \(U\) has exactly one
increasing enumeration, so this divides the \(N\)-combinations into
\(2^q\) mutually exclusive classes.  By Theorem~A
(\cref{thm:p03-w03-p03-028-ramsey-a}) there is an infinite \(Y\subseteq
\Gamma\) all of whose \(N\)-combinations lie in one \(C_\varepsilon\).  A
\(\lhd\)-increasing \(N\)-tuple \(\bar a\) from \(Y\) has distinct entries
and is the increasing enumeration of the \(N\)-combination of its entries.
So for two such tuples \(\bar a,\bar b\) both combinations lie in
\(C_\varepsilon\), and \(\psi_j(\bar a)\) and \(\psi_j(\bar b)\) both have
truth value \(\varepsilon_j\).
\end{proof}
\end{annexitem}

\begin{annexitem}{annex:p03-w03-p03-028-finite-sat}{thm:p03-w03-p03-028-finite-sat}{Finite satisfiability of the indiscernibility diagram}
\begin{proof}
Let \(\Sigma_0\subseteq\Sigma_X\) be finite.  Let \(X_0\) be the finite set
of \(x\in X\) with \(c_x\) occurring in \(\Sigma_0\), with
\(x^1<\cdots<x^p\) its elements.  Let the sentences of the third kind in
\(\Sigma_0\) be, for \(j=1,\ldots,q\),
\[
 \psi_j(c_{x_{j,1}},\ldots,c_{x_{j,n_j}})\leftrightarrow
 \psi_j(c_{y_{j,1}},\ldots,c_{y_{j,n_j}}),
\]
with \(x_{j,1}<\cdots<x_{j,n_j}\) and \(y_{j,1}<\cdots<y_{j,n_j}\) in
\(X_0\).  Let \(N=\max(1,n_1,\ldots,n_q)\) and regard each \(\psi_j\) as a
formula with free variables among \(v_0,\ldots,v_{N-1}\); its truth at an
\(N\)-tuple depends only on the first \(n_j\) entries.

Let \(\mathcal M\) be an infinite model of \(T\) and \(\lhd\) a linear
order of \(\dom\mathcal M\), which exists by the axiom of choice.
\Cref{thm:p03-w03-p03-028-ramsey-step} gives an infinite \(Y\) homogeneous
for \(\psi_1,\ldots,\psi_q\) on \(\lhd\)-increasing \(N\)-tuples.  Choose
\(y^1\lhd y^2\lhd\cdots\lhd y^{p+N}\) in \(Y\).  Interpret \(c_{x^t}\) as
\(y^t\) for \(t\leq p\), and every other new constant arbitrarily; those
constants do not occur in \(\Sigma_0\).

The sentences of \(T\) hold, and the inequalities between constants of
\(X_0\) hold because the \(y^t\) are distinct.  Fix \(j\) and write
\(x_{j,l}=x^{t_l}\) and \(y_{j,l}=x^{s_l}\), so \(t_1<\cdots<t_{n_j}\leq p\)
and \(s_1<\cdots<s_{n_j}\leq p\).  The padded tuples
\[
 \bar a=(y^{t_1},\ldots,y^{t_{n_j}},y^{p+1},\ldots,y^{p+N-n_j}),
 \qquad
 \bar b=(y^{s_1},\ldots,y^{s_{n_j}},y^{p+1},\ldots,y^{p+N-n_j})
\]
are \(\lhd\)-increasing \(N\)-tuples from \(Y\).  By homogeneity
\(\mathcal M\entails\psi_j(\bar a)\leftrightarrow\psi_j(\bar b)\), and
since \(\psi_j\) only reads the first \(n_j\) entries this is the
\(j\)-th sentence under the chosen interpretation.  So the expansion of
\(\mathcal M\) is a model of \(\Sigma_0\).
\end{proof}
\end{annexitem}

\begin{annexitem}{annex:p03-w03-p03-028-coherent}{thm:p03-w03-p03-028-coherent}{An indiscernible sequence determines a coherent blueprint}
\begin{proof}
(1) Since \(I\) is infinite it has increasing \(n\)-tuples.  If
\(\varphi\) holds of \((a_{i_1},\ldots,a_{i_n})\) for one increasing tuple,
it holds of every increasing tuple by quantifier-free order
indiscernibility.  So \(p_n\) is the quantifier-free type of
\((a_{i_1},\ldots,a_{i_n})\) for any \(i_1<\cdots<i_n\).  It is complete,
since for each quantifier-free \(\varphi\) either \(\varphi\) or
\(\neg\varphi\) holds of that tuple, and it contains \(v_s\neq v_t\) for
\(s\neq t\) because the \(a_i\) are distinct.

(2) Let \(\varphi(v_1,\ldots,v_m)\in p_m\) and put
\(\varphi'(v_1,\ldots,v_n)=\varphi(v_{k_1},\ldots,v_{k_m})\).  For
\(i_1<\cdots<i_n\), the tuple \((a_{i_{k_1}},\ldots,a_{i_{k_m}})\) is
increasing, so it satisfies \(\varphi\); hence
\(\mathcal N\entails\varphi'(a_{i_1},\ldots,a_{i_n})\) and
\(\varphi'\in p_n\).  As \((b_1,\ldots,b_n)\) realizes \(p_n\),
\(\mathcal K\entails\varphi(b_{k_1},\ldots,b_{k_m})\).  Thus
\((b_{k_1},\ldots,b_{k_m})\) satisfies every formula of the complete type
\(p_m\), so it realizes \(p_m\).
\end{proof}
\end{annexitem}

\begin{annexitem}{annex:p03-w03-p03-028-realize}{thm:p03-w03-p03-028-realize}{Realizing a coherent blueprint over any order}
\begin{proof}
Add new constants \(d_j\) (\(j\in J\)) and let \(\Theta\) consist of
\(T^{\mathrm{sk}}\) together with \(\varphi(d_{j_1},\ldots,d_{j_n})\) for
every \(n\), every \(\varphi\in p_n\) and every \(j_1<\cdots<j_n\) in
\(J\).  Let \(\Theta_0\subseteq\Theta\) be finite, mentioning the constants
for \(j^1<\cdots<j^p\).  Because \(I\) is infinite it contains
\(i^1<\cdots<i^p\).  Interpret \(d_{j^t}\) as \(a_{i^t}\) in
\(\mathcal N\).  A sentence \(\varphi(d_{j^{t_1}},\ldots,d_{j^{t_n}})\) of
\(\Theta_0\) has \(t_1<\cdots<t_n\), so \((a_{i^{t_1}},\ldots,
a_{i^{t_n}})\) is increasing and satisfies \(\varphi\in p_n\) by
\cref{thm:p03-w03-p03-028-coherent}(1).  So \(\Theta_0\) has a model, and by
compactness (\cref{thm:fol-compactness}) \(\Theta\) has one.  Let
\(\mathcal K\) be its \(\Lang^{\mathrm{sk}}\)-reduct and
\(b_j=d_j^{\mathcal K}\).  Each increasing \(n\)-tuple of the \(b_j\)
satisfies all of the complete type \(p_n\), hence realizes it; the \(b_j\)
are distinct because \(p_2\) contains \(v_1\neq v_2\).

The term hull \(H\) of \(\{b_j\}\) is the set of values of
\(\Lang^{\mathrm{sk}}\)-terms at finite tuples of the \(b_j\), including
constant terms (\cref{def:p03-morley-skolem-hull}); it contains the
interpretations of all constants and is closed under all function symbols,
so it is the universe of a substructure of \(\mathcal K\).  The axioms of
\(T^{\mathrm{sk}}\) are universal, and universal sentences true in a
structure are true in each of its substructures, so \(H\entails
T^{\mathrm{sk}}\).  Every \(\Lang\)-reduct of a model of
\(T^{\mathrm{sk}}\) is a model of \(T\)
(\cref{def:p03-morley-skolem-hull}), so the \(\Lang\)-reduct of \(H\) is a
model of \(T\).
\end{proof}
\end{annexitem}

\begin{annexitem}{annex:p03-w03-p03-028-em-57}{thm:p03-w03-p03-028-em-57}{Ehrenfeucht--Mostowski Theorem 5.7}
\begin{proof}
The language of \(S\) may have any cardinality.  Form the Skolem expansion
\(S^{\#}\) as in the proof of \cref{thm:p03-morley-ramsey-em}: put each
axiom in prenex form and replace each existentially quantified variable by
a new function symbol of the preceding universally quantified variables (a
constant if there are none), adding a fresh constant if the language has
none.  That construction does not use the cardinality of the language.
The axioms of \(S^{\#}\) are universal, every model of \(S\) expands to a
model of \(S^{\#}\) by choosing witnesses, and every reduct of a model of
\(S^{\#}\) satisfies \(S\).  In particular \(S^{\#}\) has an infinite
model.

Apply \cref{thm:p03-w03-p03-028-finite-sat} and compactness to \(S^{\#}\)
and \((X,\prec)\): there is \(\mathcal A\entails S^{\#}\) with distinct
elements \(c_x\) (\(x\in X\)), order indiscernible for all
\(\Lang^{\#}\)-formulas.  Identify \(x\) with \(c_x\), so
\(X\subseteq\dom\mathcal A\).  Let \(H\) be the term hull of \(X\); as in
the proof of \cref{thm:p03-w03-p03-028-realize}, \(H\entails S^{\#}\), and
its \(\Lang\)-reduct \(\mathcal M_0\) is a model of \(S\) with
\(X\subseteq\dom\mathcal M_0\).

Let \(g\in G(X,\prec)\).  Define \(\hat g\) on \(H\) by
\(\hat g(t(x_1,\ldots,x_n))=t(g(x_1),\ldots,g(x_n))\).  To see that this is
well defined, suppose \(t(\bar x)=t'(\bar y)\).  Let \(z_1\prec\cdots\prec
z_p\) enumerate the entries of \(\bar x\) and \(\bar y\), and write the
equation as an \(\Lang^{\#}\)-formula \(\theta(z_1,\ldots,z_p)\).  Since
\(g\) preserves \(\prec\), \(g(z_1)\prec\cdots\prec g(z_p)\), and order
indiscernibility gives \(\theta(g(z_1),\ldots,g(z_p))\), that is,
\(t(g\bar x)=t'(g\bar y)\).  The same argument with
\(t(\bar x)\neq t'(\bar y)\), and with
\(R(t_1(\bar x),\ldots,t_l(\bar x))\) or its negation for each relation
symbol \(R\), shows that \(\hat g\) is injective and preserves and reflects
every relation.  It commutes with every function symbol \(F\) by
definition, since \(F(t_1,\ldots,t_l)\) is again a term.  It is onto,
because \(t(\bar x)=\hat g(t(g^{-1}\bar x))\) and \(g^{-1}\in
G(X,\prec)\).  So \(\hat g\) is an automorphism of \(H\), hence of its
reduct \(\mathcal M_0\), and \(\hat g(x)=g(x)\) for \(x\in X\).  Thus every
element of \(G(X,\prec)\) extends to an element of
\(G_{\mathcal M_0}\).

\emph{The article's own proof} \cite[pp.~63--66]{EhrenfeuchtMostowski1956}.  By
\cref{thm:p03-w03-p03-028-em-open-reduction} it suffices to treat an open
theory \(S\).  If \(X\) is empty, any model of \(S\) serves, the only
transformation of \(\varnothing\) being extended by the identity.  Otherwise,
by Theorem~5.5 (\cref{thm:p03-w03-p03-028-em-55}) it suffices that \(S^*(X)\)
with all equivalences (5.4.1) for
\(h\in G(X,\prec)\) is consistent \cite[p.~63]{EhrenfeuchtMostowski1956}.  A derivation of
\(\bottom\) would use the closures of finitely many axioms \(\Phi\).  List the
equivalences in \(\Phi\) as \(i=1,\ldots,s\), with open matrices \(\zeta_i\),
maps \(g_i\in G(X,\prec)\) and assignments; let \(t\geq1\) bound the indices of
the variables of the \(\zeta_i\), and let \(x_{(i-1)t+1},\ldots,x_{it}\) be the
values the \(i\)-th assignment gives to \(\xi_1,\ldots,\xi_t\).  For each
element \(x\) of \(X\) occurring in \(\Phi\) but not among these values, add an
instance with \(\zeta=\xi_1\iota\xi_1\), the identity of \(X\), and all values
equal to \(x\); if no equivalence and no element of \(X\) occurs in \(\Phi\),
add one such instance for some \(x\in X\).  Then \(\Phi\) consists of axioms of
the theory \(S^*\) of \cref{thm:p03-w03-p03-028-em-57-fragment}, which has a
model; so \(\Phi\) is consistent, a contradiction.
\end{proof}
\end{annexitem}

\begin{annexitem}{annex:p03-w03-p03-028-em-56}{thm:p03-w03-p03-028-em-56}{Ehrenfeucht--Mostowski Theorem 5.6}
\begin{proof}
Apply \cref{thm:p03-w03-p03-028-em-57} with \(X=\mathbb Z\) under its
usual order.  The shift \(h(x)=x+1\) lies in \(G(\mathbb Z,<)\), and
\(h^n(0)=n\), so the powers of \(h\) are pairwise distinct and form an
infinite cyclic group \(G\subseteq G(\mathbb Z,<)\).  Every element of
\(G\) extends to an automorphism of \(\mathcal M_0\), so
\(G_{\mathcal M_0}\) strongly contains \(G\).

\emph{The article's own proof} \cite[pp.~60--62]{EhrenfeuchtMostowski1956}.  By
\cref{thm:p03-w03-p03-028-em-open-reduction} it suffices to treat an open
theory \(S\).  Let \(X=\{x_j:j\in\mathbb Z\}\) with the \(x_j\) distinct and
\(h(x_j)=x_{j+1}\), and let \(S_h\) be \(S^*(X)\) with all equivalences
(5.4.1) for this \(h\) added
\cite[p.~60]{EhrenfeuchtMostowski1956}.  \(S_h\) is consistent: a derivation of \(\bottom\) would
use the closures of finitely many axioms \(\Phi\).  Let \(t\geq1\) bound the
indices of the variables of the open matrices \(\zeta_1,\ldots,\zeta_s\) of the
equivalences in \(\Phi\), let the \(l\)-th assignment give
\(x_{i_l},x_{j_l},\ldots,x_{m_l}\) to \(\xi_1,\ldots,\xi_t\), and let \(n\) be at
least the absolute value of every index of a constant occurring in \(\Phi\).
Then \(\Phi\) consists of axioms of the theory \(S^*\) of
\cref{thm:p03-w03-p03-028-em-56-fragment}, which has a model, so \(\Phi\) is
consistent.  By \cref{thm:p03-w03-p03-028-em-complete-set} there is a
consistent and complete \(I\supseteq A(S_h)\); \(\mathfrak M(X,I)\) is a model
of \(S\) (\cref{thm:p03-w03-p03-028-em-37}), and \(h\) induces \(f_h\) with
\(f_h([x])=[h(x)]\) (\cref{thm:p03-w03-p03-028-em-54}).  For every integer
\(k\), \(f_h^k\) is an automorphism with \(f_h^k([x])=[h^k(x)]\), by induction
on \(\abs k\), using \(f_h([h^{-1}(x)])=[x]\) for negative \(k\).  The powers of
\(h\) form a cyclic group, infinite because \(h^k(x_0)=x_k\neq x_0\) for
\(k\neq0\), and each of its elements is extended by a power of \(f_h\).  After
the renaming in the proof of \cref{thm:p03-w03-p03-028-em-55}, this gives a
model \(M_0\) of \(S\) whose automorphism group strongly contains an infinite
cyclic group.  The article concludes by appeal to Theorem~5.5, whose condition
concerns every element of the group; the powers of \(f_h\) supply it.
\end{proof}
\end{annexitem}

\begin{annexitem}{annex:p03-w03-p03-028-no-finite-order}{thm:p03-w03-p03-028-no-finite-order}{Infinite cyclic cannot be replaced by an arbitrary group}
\begin{proof}
Let \(g\) be the restriction of \(f\) to \(X\), with \(g^n\) the identity
for some \(n\geq1\), and suppose \(g(x)\neq x\) for some \(x\in X\).  By
the third axiom, \(x\mathrel\pi g(x)\) or \(g(x)\mathrel\pi x\).  In the
first case apply the automorphism \(f\) repeatedly: \(g^j(x)\mathrel\pi
g^{j+1}(x)\) for every \(j\), and transitivity gives
\(x\mathrel\pi g^n(x)\), that is \(x\mathrel\pi x\), contrary to
irreflexivity.  The second case is symmetric.  So \(g\) is the identity.

If \(G_{\mathcal M}\) strongly contained a group \(G\) of transformations
of \(X\) with an element \(g\neq\mathrm{id}\) of finite order, some
automorphism \(f\) of \(\mathcal M\) would extend \(g\), and \(f\) would
map \(X\) onto itself with restriction \(g\), contradicting the first
part.  The same argument with \(\prec\) in place of \(\pi\) shows that
\(G(X,\prec)\) has no element of finite order other than the identity.
\end{proof}
\end{annexitem}

\begin{annexitem}{annex:p03-w03-p03-028-elementary}{thm:p03-w03-p03-028-elementary}{Order embeddings induce elementary embeddings}
\begin{proof}
Write \(\hat e\) for the assignment.  Let \(\theta\) be an atomic
formula of the Skolem language of \(T^{*\mathrm{sk}}\) or the negation of
one, applied to finitely many terms in the generators.  Let
\(i_1<\cdots<i_p\) enumerate all indices of generators occurring in these
terms, and write the statement as a quantifier-free formula
\(\chi(v_1,\ldots,v_p)\) evaluated at \((b_{i_1},\ldots,b_{i_p})\).  This
tuple realizes \(p_p\), so \(\chi\) holds there iff \(\chi\in p_p\).
Because \(e\) preserves order, \((b'_{e(i_1)},\ldots,b'_{e(i_p)})\) is an
increasing tuple of the second sequence and also realizes \(p_p\), so
\(\chi\) holds there iff \(\chi\in p_p\).  Hence each such statement holds
in \(H_I\) iff its image under \(\hat e\) holds in \(H_J\).

Taking \(\chi\) to be an equation between two terms shows that \(\hat e\)
is well defined, and taking its negation shows that \(\hat e\) is
injective.  Taking \(\chi\) to be an atomic relation statement or its
negation shows that \(\hat e\) preserves and reflects every relation
symbol, including every \(R_\varphi\).  It commutes with each function
symbol \(F\) by definition, since \(F\) applied to terms is again a term.
So \(\hat e\) is an embedding in the Skolem language.  If \(e\) is onto,
every term in the \(b'_j\) is the image of a term in the \(b_i\), so
\(\hat e\) is onto.  Given two hulls over the same \(J\), apply this with
the identity of \(J\).

By \cref{thm:p03-w03-p03-028-realize} applied to \(T^*\) in place of
\(T\), \(H_I\) and \(H_J\) are models of \(T^{*\mathrm{sk}}\), so their
reducts satisfy \(T^*\).  Let \(\varphi(\bar x)\) be an
\(\Lang\)-formula and \(\bar a\) a tuple from \(H_I\).  Then
\(H_I\entails\varphi(\bar a)\) iff \(H_I\entails R_\varphi(\bar a)\), by
the axioms of \(T^*\); iff \(H_J\entails R_\varphi(\hat e\bar a)\), since
\(\hat e\) preserves and reflects \(R_\varphi\); iff
\(H_J\entails\varphi(\hat e\bar a)\), again by \(T^*\).  So the
restriction to the \(\Lang\)-reducts is elementary.
\end{proof}
\end{annexitem}

The following items prove the results of
\cref{sec:p03-w03-p03-028-em1956,sec:p03-w03-p03-028-em-automorphisms,sec:p03-w03-p03-028-em-induced,sec:p03-w03-p03-028-em-original}
in the order of the article \cite[pp.~51--68]{EhrenfeuchtMostowski1956}.  The
reading of the article's syntax is that of
\cref{def:p03-w03-p03-028-em-syntax}: \(T(S)\) consists of the matrices whose
universal closures are derivable from the universal closures of the axioms.

\begin{annexitem}{annex:p03-w03-p03-028-em-21}{thm:p03-w03-p03-028-em-21}{Ehrenfeucht--Mostowski Theorem 2.1}
\begin{proof}
The proof uses only \(F(S)\subseteq F(S')\), \(P(S)\subseteq P(S')\) and
\(M'_\varphi=M_\varphi\) for \(\varphi\in F(S)\cup P(S)\); it does not use
\(T(S)\subseteq T(S')\).  First, \(\operatorname{val}_{fM'}\omega=
\operatorname{val}_{fM}\omega\) for every \(\omega\in W(S)\) and every
valuation \(f\), by induction on \(\omega\): both sides equal \(f(\xi_j)\) for
\(\omega=\xi_j\), and for \(\omega=\varphi(\omega_1,\ldots,\omega_k)\) with
\(\varphi\in F(S)\) both sides equal
\(M_\varphi(\operatorname{val}_{fM}\omega_1,\ldots,\operatorname{val}_{fM}\omega_k)\)
because \(M'_\varphi=M_\varphi\).  Next, by induction on \(\zeta\in Z(S)\),
simultaneously for all valuations: for \(\zeta=\pi(\omega_1,\ldots,\omega_k)\)
with \(\pi\in P(S)\) use \(M'_\pi=M_\pi\) and the term case; the clauses for
\(\mathord\sim\) and \(\cdot\) are the same on both sides; and
\(\operatorname{stsf}_{fM'}(\exists\xi_n)\zeta\) holds iff
\(\operatorname{stsf}_{f'M'}\zeta\) for some \(f'\sim_nf\), iff (by the
induction hypothesis at \(f'\)) \(\operatorname{stsf}_{f'M}\zeta\) for some
such \(f'\), iff \(\operatorname{stsf}_{fM}(\exists\xi_n)\zeta\); the
valuations \(f'\sim_nf\) are the same for \(M\) and \(M'\), which share the
set \(X\).
\end{proof}
\end{annexitem}

\begin{annexitem}{annex:p03-w03-p03-028-em-theorems-valid}{thm:p03-w03-p03-028-em-theorems-valid}{Theorems are the matrices valid in all models}
\begin{proof}
For a matrix \(\zeta\) whose free variables are among
\(\xi_{j_1},\ldots,\xi_{j_k}\), let \(\forall\zeta\) be its universal closure
(\(\forall\zeta=\zeta\) if \(\zeta\) is closed), and let \(\Sigma_S\) be the set
of closures of members of \(A(S)\).  By the clause for the universal quantifier
applied \(k\) times, \(M\entails\forall\zeta\) under an assignment \(s\) iff
\(M\) satisfies \(\zeta\) under every assignment that differs from \(s\) at
most at \(\xi_{j_1},\ldots,\xi_{j_k}\); since every assignment differs from
itself nowhere, \(M\entails\forall\zeta\) under every assignment iff
\(\zeta\in V_M\).  So \(M\) is a model of \(S\) iff \(M\entails\Sigma_S\).

If \(\zeta\in T(S)\), then \(\Sigma_S\derives\forall\zeta\), hence
\(\Sigma_S\entails\forall\zeta\) by soundness (\cref{thm:fol-soundness}), and
\(\zeta\in V_M\) for every model \(M\).  Conversely, if \(\zeta\in V_M\) for
every model, then every structure satisfying \(\Sigma_S\) satisfies
\(\forall\zeta\), so \(\Sigma_S\entails\forall\zeta\), and
\(\Sigma_S\derives\forall\zeta\) by completeness
(\cref{thm:fol-completeness}, which assumes the axiom of choice).  If
\(S\) has a model \(M\) and \(\Sigma_S\derives\bottom\), soundness gives
\(M\entails\bottom\), which is impossible; so \(S\) is consistent.  If \(S\) has
no model, \(\Sigma_S\entails\bottom\) holds vacuously and completeness gives
\(\Sigma_S\derives\bottom\).
\end{proof}
\end{annexitem}

\begin{annexitem}{annex:p03-w03-p03-028-em-11}{thm:p03-w03-p03-028-em-11}{Ehrenfeucht--Mostowski Theorem 1.1}
\begin{proof}
For each \(\zeta\in A(S)\) let \(\sigma_\zeta\) be a rectified prenex sentence
logically equivalent to the universal closure of \(\zeta\)
(\cref{thm:prenex-normal-form}), and let \(\operatorname{Sk}(\sigma_\zeta)\) be
its Skolemization (\cref{thm:skolem-equisatisfiable}), the new functors being
chosen so that different axioms receive disjoint sets of new functors, none
in \(F(S)\).  \(\operatorname{Sk}(\sigma_\zeta)\) has the form
\(\forall u_1\cdots\forall u_k\,\chi_\zeta\) with \(\chi_\zeta\)
quantifier-free; renaming \(u_1,\ldots,u_k\) as variables \(\xi_j\), we take
\(\chi_\zeta\) to be an open matrix.  Let \(S'\) have
\(F(S')=F(S)\cup\{\text{new functors}\}\), \(P(S')=P(S)\) and
\(A(S')=\{\chi_\zeta:\zeta\in A(S)\}\).  \(S'\) is open.  For a pseudo-model
\(N\) of \(S'\), \(\chi_\zeta\in V_N\) iff \(N\entails\operatorname{Sk}(\sigma_\zeta)\),
by the first paragraph of the proof of
\cref{thm:p03-w03-p03-028-em-theorems-valid}.

(a) \emph{Every model \(M\) of \(S\) over a set \(Y\) is the restriction of a
model of \(S'\) over \(Y\).}  \(M\entails\sigma_\zeta\) for each axiom
\(\zeta\).  By \cref{thm:skolem-equisatisfiable}(1), for each \(\zeta\)
there is an interpretation of the new functors of \(\sigma_\zeta\) under which
\(\operatorname{Sk}(\sigma_\zeta)\) holds; by the axiom of choice choose one
for every \(\zeta\) at once.  As the sets of new functors are disjoint, these
choices together with \(M\) define a pseudo-model \(M'\) of \(S'\).  For each
\(\zeta\), \(M'\) agrees on the symbols of
\(\operatorname{Sk}(\sigma_\zeta)\) with the expansion chosen for \(\zeta\),
so by \cref{thm:p03-w03-p03-028-em-21} (which uses only the inclusion of
symbols) \(\chi_\zeta\in V_{M'}\).  So \(M'\) is a model of \(S'\).

(b) \emph{\(T(S)\subseteq T(S')\).}  Let \(N\) be a model of \(S'\) and \(N_0\)
its restriction to \(F(S)\cup P(S)\).  By
\cref{thm:skolem-equisatisfiable}(2) and
\cref{thm:p03-w03-p03-028-em-21}, \(N_0\entails\sigma_\zeta\), hence \(N_0\)
satisfies the closure of \(\zeta\), for each \(\zeta\in A(S)\); so \(N_0\) is a
model of \(S\).  If \(\zeta\in T(S)\), then \(\zeta\in V_{N_0}\) by
\cref{thm:p03-w03-p03-028-em-theorems-valid}, hence \(\zeta\in V_N\) by
\cref{thm:p03-w03-p03-028-em-21}.  As \(N\) was arbitrary, \(\zeta\in T(S')\)
by \cref{thm:p03-w03-p03-028-em-theorems-valid}.  With \(F(S)\subseteq F(S')\)
and \(P(S)=P(S')\), \(S'\) is an extension of \(S\).

(c) \emph{\(T(S')\cap Z(S)\subseteq T(S)\).}  Let \(\zeta\in Z(S)\cap T(S')\)
and let \(M\) be any model of \(S\).  By (a) \(M\) is the restriction of a
model \(M'\) of \(S'\); \(\zeta\in V_{M'}\) by
\cref{thm:p03-w03-p03-028-em-theorems-valid}, so \(\zeta\in V_M\) by
\cref{thm:p03-w03-p03-028-em-21}.  Hence \(\zeta\in T(S)\).  With (b),
\(T(S')\cap Z(S)=T(S)\), and the extension is inessential.
\end{proof}
\end{annexitem}

\begin{annexitem}{annex:p03-w03-p03-028-em-31}{thm:p03-w03-p03-028-em-31}{Ehrenfeucht--Mostowski Lemma 3.1}
\begin{proof}
Let \(M\) be a model of \(S\) over an infinite set \(Y\).  First let
\(X=\{x_1,\ldots,x_n\}\) be finite, and choose distinct \(y_1,\ldots,y_n\in Y\).
Extend \(M\) to a pseudo-model \(M^*\) of \(S^*(X)\) over \(Y\) by
\(M^*_\varphi=M_\varphi\) for \(\varphi\in F(S)\), \(M^*_{x_j}=y_j\), and
\(M^*_\pi=M_\pi\) for \(\pi\in P(S)\).  By
\cref{thm:p03-w03-p03-028-em-21}, each \(\zeta\in A(S)\) lies in \(V_{M^*}\);
the matrices \(\mathord\sim(x'\iota x'')\) lie in \(V_{M^*}\) because the
\(y_j\) are distinct.  So \(M^*\) is a model of \(S^*(X)\), and \(S^*(X)\) is
consistent by \cref{thm:p03-w03-p03-028-em-theorems-valid} (only its
soundness direction, which uses no choice, is needed).

In general, suppose the closures of the axioms of \(S^*(X)\) derive
\(\bottom\).  The derivation has finitely many open assumptions; let \(X_1\) be
the finite set of elements of \(X\) occurring in them.  These assumptions are
closures of axioms of \(S^*(X_1)\), and the language of the derivation extends
that of \(S^*(X_1)\) only by constants, so by
\cref{lem:new-constant-elimination} \(S^*(X_1)\) derives \(\bottom\) in
its own language.  This contradicts the finite case.
\end{proof}
\end{annexitem}

\begin{annexitem}{annex:p03-w03-p03-028-em-complete-set}{thm:p03-w03-p03-028-em-complete-set}{Complete sets exist}
\begin{proof}
Let \(\Sigma\) be the set of universal closures of members of \(A(S_1)\); it is
consistent.  By the axiom of choice the formulas of the language of \(S_1\)
can be enumerated by an ordinal, and \cref{lem:transfinite-lindenbaum}
gives a maximal consistent \(\Delta\supseteq\Sigma\).  Put
\(I=\{\zeta\in Z(S_1):\forall\zeta\in\Delta\}\).  Then \(A(S_1)\subseteq I\).
The closures of members of \(I\) lie in \(\Delta\), so the theory with axioms
\(I\) is consistent.  If \(\zeta\) is a theorem of that theory, \(\forall\zeta\)
is derivable from members of \(\Delta\).  The closure properties of
\cref{lem:prop-mcs} hold for first-order maximal consistent sets
(\cref{sec:fol-metatheory-01}); by clause~1, \(\forall\zeta\in\Delta\), and
\(\zeta\in I\).  For closed \(\zeta\), \(\forall\zeta=\zeta\), and by clause~2
one of \(\zeta,\neg\zeta\) lies in \(\Delta\), so \(\zeta\in I\) or \(\mathord\sim\zeta\in I\).  For the last
sentence take \(S_1=S^*(X)\), which is consistent by
\cref{thm:p03-w03-p03-028-em-31}.
\end{proof}
\end{annexitem}

\begin{annexitem}{annex:p03-w03-p03-028-em-32}{thm:p03-w03-p03-028-em-32}{Ehrenfeucht--Mostowski Lemmas 3.2 and 3.3}
\begin{proof}
We first record five properties of any consistent and complete \(I\) in the
sense of \cref{def:p03-w03-p03-028-em-mxi}; they are used here and in the
following items.  Let \(\sigma,\sigma_1,\sigma_2\) be closed matrices of
\(S^*(X)\).
\begin{enumerate}
\item[(I1)] If \(\zeta\in I\) and \(\tau_1,\tau_2,\ldots\) are constant terms,
then \(\operatorname{subst}\zeta(\tau_1,\tau_2,\ldots)\in I\): it is derived
from \(\forall\zeta\) by \(\forall\)-elimination, the constant terms being
free for the variables, so it is a theorem of \(S(X,I)\).
\item[(I2)] \(\mathord\sim\sigma\in I\) iff \(\sigma\notin I\): at least one of
the two is in \(I\) by completeness, and not both, since together they derive
\(\bottom\).
\item[(I3)] \(\sigma_1\cdot\sigma_2\in I\) iff \(\sigma_1\in I\) and
\(\sigma_2\in I\): each direction is a derivation, and \(I\) contains its
theorems.
\item[(I4)] If \(\sigma_1\in I\) and \(\sigma_1\supset\sigma_2\in I\), then
\(\sigma_2\in I\).
\item[(I5)] \(\sigma_1\equiv\sigma_2\in I\) iff (\(\sigma_1\in I\) iff
\(\sigma_2\in I\)).  If the equivalence is in \(I\) and one side is, the other
is derivable.  If both sides are in \(I\), the equivalence is derivable; if
neither is, \(\mathord\sim\sigma_1,\mathord\sim\sigma_2\in I\) by (I2), and
again the equivalence is derivable.
\end{enumerate}

\emph{Lemma~3.2.}  The five identity matrices of
\cref{def:p03-w03-p03-028-em-syntax} are theorems of \(S(X,I)\) and so lie in
\(I\).  From \(\xi_1\iota\xi_1\in I\) and (I1), \(\omega\iota\omega\in I\) for
every constant term \(\omega\).  From
\(\xi_1\iota\xi_2\supset\xi_2\iota\xi_1\in I\), (I1) and (I4),
\(\omega_1\approx\omega_2\) implies \(\omega_2\approx\omega_1\); in the same
way the third identity matrix gives transitivity.  For distinct
\(x_1,x_2\in X\), \(\mathord\sim(x_1\iota x_2)\in A(S^*(X))\subseteq I\), so
\(x_1\iota x_2\notin I\) by (I2), that is, \(x_1\not\approx x_2\).

\emph{Lemma~3.3.}  Let \(k=a(\varphi)\) and let \(\omega_j\approx\omega'_j\)
for \(j\leq k\).  For \(j\leq k\) let \(\rho_j\) be the term
\(\varphi(\omega'_1,\ldots,\omega'_{j-1},\omega_j,\ldots,\omega_k)\), so
\(\rho_1=\varphi(\omega_1,\ldots,\omega_k)\), and let
\(\rho_{k+1}=\varphi(\omega'_1,\ldots,\omega'_k)\).  The fourth identity
matrix with \(\omega=\varphi(\xi_1,\ldots,\xi_k)\), \(\xi_j\) in the role of
\(\xi_k\) and a variable \(\xi_l\) not among \(\xi_1,\ldots,\xi_k\) in the role of
\(\xi_l\), is in \(I\); substituting \(\omega_j\) for \(\xi_j\), \(\omega'_j\) for
\(\xi_l\), \(\omega'_i\) for \(\xi_i\) (\(i<j\)) and \(\omega_i\) for \(\xi_i\)
(\(i>j\)) gives, by (I1), \(\omega_j\iota\omega'_j\supset\rho_j\iota\rho_{j+1}\in
I\), so \(\rho_j\approx\rho_{j+1}\) by (I4).  By transitivity
\(\rho_1\approx\rho_{k+1}\).  For \(\pi\) with \(a(\pi)=k\) and
\(\tau_j\approx\tau'_j\), the fifth identity matrix with
\(\zeta=\pi(\xi_1,\ldots,\xi_k)\) gives in the same way, for each \(j\),
\(\pi(\tau'_1,\ldots,\tau'_{j-1},\tau_j,\ldots,\tau_k)\equiv
\pi(\tau'_1,\ldots,\tau'_j,\tau_{j+1},\ldots,\tau_k)\in I\) by (I1) and (I4), and
(I5) turns the chain of these equivalences into
\(\pi(\tau_1,\ldots,\tau_k)\in I\equiv\pi(\tau'_1,\ldots,\tau'_k)\in I\).
\end{proof}
\end{annexitem}

\begin{annexitem}{annex:p03-w03-p03-028-em-34}{thm:p03-w03-p03-028-em-34}{Ehrenfeucht--Mostowski Lemmas 3.4--3.6}
\begin{proof}
\emph{Lemma~3.4.}  \(\mathfrak X_X\) is non-empty by the side condition of
\cref{def:p03-w03-p03-028-em-mxi}.  By
\cref{thm:p03-w03-p03-028-em-32}, the values of \(\mathfrak M_\varphi\) and
\(\mathfrak M_\pi\) do not depend on the representatives \(\omega_j\) chosen in
the classes \([\omega_j]\), so both are well defined.  For the identity
predicate, \(\mathfrak M_\iota([\omega_1],[\omega_2])\) holds iff
\(\omega_1\iota\omega_2\in I\), iff \(\omega_1\approx\omega_2\), iff
\([\omega_1]=[\omega_2]\) \cite[pp.~54--55]{EhrenfeuchtMostowski1956}.  So
\(\mathfrak M\) assigns the identity of \(\mathfrak X_X\) to \(\iota\) and is a
pseudo-model of \(S\).

\emph{Lemma~3.5}, by induction on \(\omega\in W(S)\).  For \(\omega=\xi_j\)
both sides are \([\tau_j]\).  For \(\omega=\varphi(\omega_1,\ldots,\omega_k)\),
\(\operatorname{val}_{\mathfrak M}\omega([\tau_1],\ldots)
=\mathfrak M_\varphi(\operatorname{val}_{\mathfrak M}\omega_1([\tau_1],\ldots),\ldots)
=\mathfrak M_\varphi([\operatorname{subst}\omega_1(\tau_1,\ldots)],\ldots)
=[\varphi(\operatorname{subst}\omega_1(\tau_1,\ldots),\ldots)]
=[\operatorname{subst}\omega(\tau_1,\ldots)]\).

\emph{Lemma~3.6}, by induction on the open matrix \(\zeta\).  Since \(\zeta\)
has no bound variables and every variable is replaced by a constant term,
\(\operatorname{subst}\zeta(\tau_1,\tau_2,\ldots)\) is closed.  For
\(\zeta=\pi(\omega_1,\ldots,\omega_k)\), Lemma~3.5 gives
\(\operatorname{stsf}_{\mathfrak M}\zeta([\tau_1],\ldots)\equiv
\mathfrak M_\pi([\operatorname{subst}\omega_1(\tau_1,\ldots)],\ldots)\equiv
\pi(\operatorname{subst}\omega_1(\tau_1,\ldots),\ldots)\in I\), and the last
matrix is \(\operatorname{subst}\zeta(\tau_1,\ldots)\).  For
\(\mathord\sim\zeta\) use (I2) of the proof of
\cref{thm:p03-w03-p03-028-em-32}, and for \(\zeta_1\cdot\zeta_2\) use (I3).
\end{proof}
\end{annexitem}

\begin{annexitem}{annex:p03-w03-p03-028-em-37}{thm:p03-w03-p03-028-em-37}{Ehrenfeucht--Mostowski Theorem 3.7}
\begin{proof}
\(\mathfrak M(X,I)\) is a pseudo-model of \(S\) by
\cref{thm:p03-w03-p03-028-em-34}.  Let \(\zeta\in A(S)\), an open matrix, and
let \(f\) be a valuation in \(\mathfrak X_X\).  Choose \(\tau_j\in f(\xi_j)\)
for every \(j\), by the axiom of choice.  By Lemma~3.6,
\(\operatorname{stsf}_{f\mathfrak M}\zeta\equiv
\operatorname{subst}\zeta(\tau_1,\tau_2,\ldots)\in I\), and the right side
holds by (I1), because \(\zeta\in A(S)\subseteq A(S^*(X))\subseteq I\).  So
\(A(S)\subseteq V_{\mathfrak M}\) \cite[p.~55]{EhrenfeuchtMostowski1956}.
\end{proof}
\end{annexitem}

\begin{annexitem}{annex:p03-w03-p03-028-em-38}{thm:p03-w03-p03-028-em-38}{Ehrenfeucht--Mostowski Theorem 3.8}
\begin{proof}
Every variable of \(\omega\) is replaced by an element of \(X\), which is a
constant of \(S^*(X)\), so \(\omega'\) is constant; as \(\zeta\) is open, the
same holds for \(\zeta'\), which is closed.  The value of a constant term and
the satisfaction of a closed open matrix do not depend on the valuation, by
the same induction.  For (3.8.1) argue by induction on \(\omega\).  If
\(\omega=\xi_j\), then \(\omega'=x_j\) and both sides equal \(M'_{x_j}\).  If
\(\omega=\varphi(\omega_1,\ldots,\omega_k)\), then
\(\omega'=\varphi(\omega'_1,\ldots,\omega'_k)\), the accents denoting the same
substitution, and if (3.8.1) holds for the \(\omega_j\),
\[
  \operatorname{val}_{M'}\omega'
  =M'_\varphi(\operatorname{val}_{M'}\omega'_1,\ldots)
  =M'_\varphi\bigl(\operatorname{val}_M\omega_1(M'_{x_1},M'_{x_2},\ldots),\ldots\bigr),
\]
which is \(\operatorname{val}_M\omega(M'_{x_1},M'_{x_2},\ldots)\) because
\(M'_\varphi=M_\varphi\) \cite[p.~55]{EhrenfeuchtMostowski1956}.  For (3.8.2)
argue by induction on the open matrix \(\zeta\): for
\(\zeta=\pi(\omega_1,\ldots,\omega_k)\) use (3.8.1) and \(M'_\pi=M_\pi\); the
clauses for \(\mathord\sim\) and \(\cdot\) are the same on both sides.
\end{proof}
\end{annexitem}

\begin{annexitem}{annex:p03-w03-p03-028-em-41}{thm:p03-w03-p03-028-em-41}{Ehrenfeucht--Mostowski Lemmas 4.1 and 4.2}
\begin{proof}
\emph{Lemma~4.1.}  For terms, induction on \(\omega\): for \(\omega=\xi_j\) both
sides are \(f(x_j)\); for \(\omega=\varphi(\omega_1,\ldots,\omega_k)\),
\begin{align*}
  f\bigl(M_\varphi(\operatorname{val}_M\omega_1(x_1,\ldots),\ldots)\bigr)
  &=M_\varphi\bigl(f(\operatorname{val}_M\omega_1(x_1,\ldots)),\ldots\bigr)\\
  &=M_\varphi\bigl(\operatorname{val}_M\omega_1(f(x_1),\ldots),\ldots\bigr).
\end{align*}
For matrices,
induction on \(\zeta\), for all valuations at once.  Atomic matrices: \(M_\pi\)
holds of the values at \(x_1,x_2,\ldots\) iff it holds of their images under
\(f\), which by the term case are the values at \(f(x_1),f(x_2),\ldots\).  The
clauses for \(\mathord\sim\) and \(\cdot\) are immediate.  For
\((\exists\xi_n)\zeta\) and a valuation \(g\): if \(g'\sim_ng\) satisfies
\(\zeta\), then \(f\circ g'\sim_nf\circ g\) satisfies \(\zeta\) by the induction
hypothesis; conversely, if \(k\sim_nf\circ g\) satisfies \(\zeta\), then
\(g'=f^{-1}\circ k\sim_ng\), and \(g'\) satisfies \(\zeta\) by the induction
hypothesis applied to \(g'\), since \(f\circ g'=k\).

\emph{Lemma~4.2.}  An \(f\in G_{M'}\) is a one-one map of \(X\) onto itself
preserving \(M'_\varphi\) and \(M'_\pi\) for all \(\varphi\in F(S')\),
\(\pi\in P(S')\).  Since \(F(S)\subseteq F(S')\), \(P(S)\subseteq P(S')\) and
\(M'\) agrees with \(M\) on these symbols, \(f\) preserves every \(M_\varphi\)
and \(M_\pi\), so \(f\in G_M\).
\end{proof}
\end{annexitem}

\begin{annexitem}{annex:p03-w03-p03-028-em-43}{thm:p03-w03-p03-028-em-43}{Ehrenfeucht--Mostowski Theorem 4.3}
\begin{proof}
For \(g_0\in G\) let \(l_{g_0}(g)=g_0g\), a left translation.  Let
\(F(S)=\varnothing\), let \(P(S)\) consist of \(\iota\) and a binary predicate
\(\pi_f\) for each one-one map \(f\) of \(G\) onto itself that is not a left
translation, and let \(A(S)\) consist of the identity matrices of
\cref{def:p03-w03-p03-028-em-syntax} \cite[p.~56]{EhrenfeuchtMostowski1956}.
If \(f\) is not a left translation, there are \(g_1,g_2\) with
\(f(g_1)g_1^{-1}\neq f(g_2)g_2^{-1}\): otherwise \(f(g)g^{-1}=g_0\) for a fixed
\(g_0\) and all \(g\), and \(f=l_{g_0}\).  By the axiom of choice select such a
pair \(g_{1f},g_{2f}\) for every \(f\), and define \(M_{\pi_f}\) by (4.3.1):
\(M_{\pi_f}(g',g'')\) iff \(g'=gg_{1f}\) and \(g''=gg_{2f}\) for some
\(g\in G\).  With \(M_\iota\) the identity of \(G\), \(M\) is a pseudo-model of
\(S\) over \(G\); the identity matrices are valid in every pseudo-model, since
\(\iota\) denotes identity (for the fifth, use
\cref{thm:fol-substitution}), so \(M\) is a model of \(S\).

\emph{Every left translation is an automorphism.}  For \(l=l_{g_0}\) and every
\(\pi_f\), \(M_{\pi_f}(l(g'),l(g''))\) holds iff
\(g_0g'=hg_{1f}\) and \(g_0g''=hg_{2f}\) for some \(h\), iff (putting
\(h=g_0g\)) \(g'=gg_{1f}\) and \(g''=gg_{2f}\) for some \(g\), iff
\(M_{\pi_f}(g',g'')\); this is (4.3.2) \cite[p.~56]{EhrenfeuchtMostowski1956}.
\(l\) is one-one and onto and preserves identity, so \(l\in G_M\).

\emph{No other map is an automorphism.}  Let \(l\) be one-one and onto but not
a left translation, and suppose \(l\in G_M\).  \(M_{\pi_l}(g_{1l},g_{2l})\) holds
(take \(g\) the unit), so \(M_{\pi_l}(l(g_{1l}),l(g_{2l}))\) holds: there is
\(g\) with \(l(g_{1l})=gg_{1l}\) and \(l(g_{2l})=gg_{2l}\), whence
\(l(g_{1l})g_{1l}^{-1}=g=l(g_{2l})g_{2l}^{-1}\), against the choice of the pair
\cite[p.~56]{EhrenfeuchtMostowski1956}.

So \(G_M\) is the group of left translations
\cite[p.~57]{EhrenfeuchtMostowski1956}.  The map \(g_0\mapsto l_{g_0}\) is a
homomorphism, since \(l_{g_0g_1}(g)=g_0g_1g=l_{g_0}(l_{g_1}(g))\); it is
one-one, since \(l_{g_0}\) sends the unit to \(g_0\); and it is onto
\(G_M\).  Hence \(G_M\) and \(G\) are isomorphic.
\end{proof}
\end{annexitem}

\begin{annexitem}{annex:p03-w03-p03-028-em-51}{thm:p03-w03-p03-028-em-51}{Ehrenfeucht--Mostowski Theorem 5.1}
\begin{proof}
The members of \(Y^n\) are the \(n\)-combinations of \(Y\), and the partition
divides them into \(k\) mutually exclusive classes \(C_1,\ldots,C_k\)
(\(k\geq1\), as \(Y^n\) is non-empty).  Ramsey's Theorem~A
(\cref{thm:p03-w03-p03-028-ramsey-a}) with \(\Gamma=Y\), \(r=n\) and \(\mu=k\),
under the axiom of choice it assumes, gives an infinite \(\Delta\subseteq Y\) all
of whose \(n\)-combinations lie in one \(C_j\).  Put \(Y_1=\Delta\); then
\(Y_1^n\subseteq C_j\).
\end{proof}
\end{annexitem}

\begin{annexitem}{annex:p03-w03-p03-028-em-52}{thm:p03-w03-p03-028-em-52}{Ehrenfeucht--Mostowski Lemmas 5.2 and 5.3}
\begin{proof}
\emph{Lemma~5.2.}  Let \(\tau\) be a constant term of \(S^*(X)\).  Its functors
lie in \(F(S)\cup X\).  Let \(x_1,\ldots,x_k\) be the distinct elements of \(X\)
occurring in \(\tau\), and let \(\omega\in W(S)\) result from \(\tau\) by
replacing each occurrence of \(x_j\) by \(\xi_j\); then
\(\tau=\operatorname{subst}\omega(x_1,\ldots,x_k,x_1,x_1,\ldots)\), the
variables \(\xi_j\) with \(j>k\) not occurring in \(\omega\) (if \(k=0\), any
elements may be substituted, or none if \(X\) is empty).  By Lemma~3.5
(\cref{thm:p03-w03-p03-028-em-34}) with \(\tau_j=x_j\),
\([\tau]=\operatorname{val}_{\mathfrak M}\omega([x_1],[x_2],\ldots)\).  If \(f\)
is an automorphism with \(f([x])=[h(x)]\), Lemma~4.1
(\cref{thm:p03-w03-p03-028-em-41}) gives
\(f([\tau])=\operatorname{val}_{\mathfrak M}\omega([h(x_1)],[h(x_2)],\ldots)\)
\cite[p.~57]{EhrenfeuchtMostowski1956}.  This value depends only on \(h\), so
two automorphisms with the property agree on every class \([\tau]\) and are
equal.

\emph{Lemma~5.3.}  If \(h_1(x)\neq h_2(x)\), then
\([h_1(x)]\neq[h_2(x)]\) by Lemma~3.2 (\cref{thm:p03-w03-p03-028-em-32}), so
\(f_{h_1}([x])\neq f_{h_2}([x])\).
\end{proof}
\end{annexitem}

\begin{annexitem}{annex:p03-w03-p03-028-em-54}{thm:p03-w03-p03-028-em-54}{Ehrenfeucht--Mostowski Lemma 5.4}
\begin{proof}
The properties (I1)--(I5) of \(I\) are those of the proof of
\cref{thm:p03-w03-p03-028-em-32}.

\emph{Necessity.}  Let \(h\) induce \(f_h\), let \(\zeta\) be open and
\(x_1,x_2,\ldots\in X\), and put
\(\sigma_1=\operatorname{subst}\zeta(x_1,x_2,\ldots)\),
\(\sigma_2=\operatorname{subst}\zeta(h(x_1),h(x_2),\ldots)\), both closed.  By
Lemma~3.6, \(\sigma_1\in I\) iff
\(\operatorname{stsf}_{\mathfrak M}\zeta([x_1],[x_2],\ldots)\); by Lemma~4.1
applied to \(f_h\), iff
\(\operatorname{stsf}_{\mathfrak M}\zeta([h(x_1)],[h(x_2)],\ldots)\); by
Lemma~3.6 again, iff \(\sigma_2\in I\).  By (I5), \(\sigma_1\equiv\sigma_2\in
I\).  (The article reaches the same conclusion by assuming, as it may by the
completeness of \(I\), that \(\sigma_1\in I\) \cite[p.~57]{EhrenfeuchtMostowski1956}.)

\emph{Sufficiency.}  Assume (5.4.1) for every open
\(\zeta\) and every assignment.  For a closed matrix or a constant term
\(\alpha\) of \(S^*(X)\) without quantifiers, let \(\alpha^h\) result from
\(\alpha\) by replacing each element \(x\) of \(X\) occurring in it by \(h(x)\).

\((\ast)\) \emph{For every closed matrix \(\sigma\) of \(S^*(X)\) without
quantifiers, \(\sigma\in I\) iff \(\sigma^h\in I\).}  Let \(x_1,\ldots,x_k\) be
the distinct elements of \(X\) in \(\sigma\), and let \(\zeta\) result from
\(\sigma\) by replacing \(x_j\) by \(\xi_j\).  \(\zeta\) is an open matrix of
\(S\), \(\sigma=\operatorname{subst}\zeta(x_1,\ldots,x_k,x_1,\ldots)\) and
\(\sigma^h=\operatorname{subst}\zeta(h(x_1),\ldots,h(x_k),h(x_1),\ldots)\), the
variables \(\xi_j\) for \(j>k\) not occurring in \(\zeta\) (if \(k=0\),
\(\sigma^h=\sigma\)).  The equivalence \(\sigma\equiv\sigma^h\) is an instance of
(5.4.1), so it is in \(I\), and (I5) gives \((\ast)\).
This is the content of (5.4.2) and of the injectivity argument
\cite[pp.~58--59]{EhrenfeuchtMostowski1956}, where the replacement of elements
by variables is written \(\tau\mapsto\bar\tau\).

Define \(f_h([\tau])=[\tau^h]\) for constant terms \(\tau\) of \(S^*(X)\)
\cite[p.~58]{EhrenfeuchtMostowski1956}.  It is well defined and one-one: by
\((\ast)\) with \(\sigma=\tau_1\iota\tau_2\), \(\tau_1\approx\tau_2\) iff
\(\tau_1^h\approx\tau_2^h\).  It is onto: replacing each \(x\) by
\(h^{-1}(x)\) and then by \(h(h^{-1}(x))=x\) returns \(\tau\), so
\([\tau]=f_h([\tau'])\) for the term \(\tau'\) obtained by the first
replacement.  \(f_h([x])=[h(x)]\).  For \(\varphi\in F(S)\),
\(f_h(\mathfrak M_\varphi([\tau_1],\ldots))=[\varphi(\tau_1,\ldots)^h]
=[\varphi(\tau_1^h,\ldots)]=\mathfrak M_\varphi(f_h([\tau_1]),\ldots)\).  For
\(\pi\in P(S)\), \(\mathfrak M_\pi([\tau_1],\ldots)\) iff
\(\pi(\tau_1,\ldots)\in I\), iff \(\pi(\tau_1^h,\ldots)\in I\) by \((\ast)\), iff
\(\mathfrak M_\pi(f_h([\tau_1]),\ldots)\) \cite[p.~59]{EhrenfeuchtMostowski1956}.
So \(f_h\) is an automorphism of \(\mathfrak M(X,I)\) with
\(f_h([x])=[h(x)]\), and \(h\) induces an automorphism.  By Lemma~5.2 the
induced automorphism is this \(f_h\), which proves the last sentence.
\end{proof}
\end{annexitem}

\begin{annexitem}{annex:p03-w03-p03-028-em-55}{thm:p03-w03-p03-028-em-55}{Ehrenfeucht--Mostowski Theorem 5.5}
\begin{proof}
Let \(S_G\) be \(S^*(X)\) with all equivalences
(5.4.1) for \(h\in G\) added to its axioms.

\emph{Sufficiency} \cite[p.~60]{EhrenfeuchtMostowski1956}.  Assume \(S_G\) is
consistent.  If \(X\) is empty and \(F(S)\) contains no constant, \(G\)
consists of the empty map, \(S_G\) is \(S\), and any model of \(S\)
(\cref{thm:p03-w03-p03-028-em-theorems-valid}) has an identity automorphism
extending it.  Otherwise \cref{thm:p03-w03-p03-028-em-complete-set} gives
\(I\supseteq A(S_G)\) that is consistent and complete in the sense of
\cref{def:p03-w03-p03-028-em-mxi}, since \(S_G\) and \(S^*(X)\) have the same
symbols.  (In section~3 the hypothesis that \(S\) has a model over an infinite
set serves only, through Lemma~3.1, to provide such an \(I\).)  By
\cref{thm:p03-w03-p03-028-em-37}, \(\mathfrak M(X,I)\) is a model of \(S\), and
since \(I\) contains every equivalence (5.4.1) for
\(h\in G\), \cref{thm:p03-w03-p03-028-em-54} shows that each \(h\in G\) induces
\(f_h\) with \(f_h([x])=[h(x)]\).

\emph{Renaming.}  By Lemma~3.2, \(x\mapsto[x]\) is one-one.  Let
\(R=(\mathfrak X_X\setminus[X])\times\{X\}\), where \([X]=\{[x]:x\in X\}\).
No \((c,X)\) lies in \(X\): otherwise \(X\in\{c,X\}\in(c,X)\in X\), a membership
cycle excluded by the axiom of foundation.  Let \(X_1=X\cup R\) and let
\(\beta\colon\mathfrak X_X\to X_1\) send \([x]\) to \(x\) and \(c\notin[X]\) to
\((c,X)\); \(\beta\) is a bijection.  Transport \(\mathfrak M(X,I)\) along
\(\beta\): \(M_\varphi=\beta\circ\mathfrak M_\varphi\circ\beta^{-1}\) on tuples,
and \(M_\pi(\beta a_1,\ldots)\) iff \(\mathfrak M_\pi(a_1,\ldots)\).  The
induction in the proof of Lemma~4.1, with \(\beta\) in place of \(f\) and two
pseudo-models in place of one, shows
\(\operatorname{stsf}_M\zeta(\beta a_1,\beta a_2,\ldots)\equiv
\operatorname{stsf}_{\mathfrak M}\zeta(a_1,a_2,\ldots)\), so \(M\) is a model of
\(S\) over \(X_1\supseteq X\).  For \(h\in G\),
\(\beta f_h\beta^{-1}\in G_M\), and \(\beta f_h\beta^{-1}(x)=\beta([h(x)])=h(x)\).
So \(G_M\) strongly contains \(G\).

\emph{Necessity} \cite[p.~60]{EhrenfeuchtMostowski1956}.  Let \(M\) be a model
of \(S\) over \(X_1\supseteq X\) with \(G_M\) strongly containing \(G\).  Extend
\(M\) to a pseudo-model \(M'\) of \(S^*(X)\) over \(X_1\) with \(M'_x=x\).  The
axioms of \(S\) are valid in \(M'\) by \cref{thm:p03-w03-p03-028-em-21}, and
\(\mathord\sim(x'\iota x'')\) is valid for \(x'\neq x''\).  For \(h\in G\) let
\(f\in G_M\) extend \(h\).  For an open \(\zeta\) and \(x_1,x_2,\ldots\in X\),
Theorem~3.8 and Lemma~4.1 give
\begin{align*}
  \operatorname{stsf}_{M'}\operatorname{subst}\zeta(x_1,x_2,\ldots)
  &\equiv\operatorname{stsf}_M\zeta(x_1,x_2,\ldots)
  \equiv\operatorname{stsf}_M\zeta(f(x_1),f(x_2),\ldots)\\
  &\equiv\operatorname{stsf}_{M'}\operatorname{subst}\zeta(h(x_1),h(x_2),\ldots),
\end{align*}
so every equivalence (5.4.1) is valid in \(M'\).  Thus
\(M'\) is a model of \(S_G\), which is consistent by
\cref{thm:p03-w03-p03-028-em-theorems-valid}.
\end{proof}
\end{annexitem}

\begin{annexitem}{annex:p03-w03-p03-028-em-56-fragment}{thm:p03-w03-p03-028-em-56-fragment}{The finite step in the proof of Theorem 5.6}
\begin{proof}
Let \(M\) be a model of \(S\) over an infinite set \(Y\), and order \(Y\) by a
linear order \(\ll\) (axiom of choice); \(\ll\) need have nothing to do with the
relations definable in \(S\) \cite[p.~61]{EhrenfeuchtMostowski1956}.  For
\(\zeta\) a matrix let \(\zeta^{+1}=\zeta\) and
\(\zeta^{-1}=\mathord\sim\zeta\).  For \(\varepsilon=(\varepsilon_1,\ldots,
\varepsilon_s)\in\{+1,-1\}^s\) let \(C_\varepsilon\) be the set of
\((2n+1)\)-element sets \(U\subseteq Y\) whose increasing enumeration
\(y_{-n}\ll\cdots\ll y_n\) satisfies
\(\operatorname{stsf}_M\zeta_l^{\varepsilon_l}(y_{i_l},y_{j_l},\ldots,y_{m_l})\)
for \(l=1,\ldots,s\).  The article expresses this condition as satisfaction of
the conjunction \(\psi_{\varepsilon_1\cdots\varepsilon_s}\) of (5.6.3) at
\((y_{-n},\ldots,y_n)\), with variables \(\bar x_p\) assigned to the indices
\(p\) \cite[(5.6.3), (5.6.6), p.~61]{EhrenfeuchtMostowski1956}; the two
formulations agree by the variable form of the induction proving Theorem~3.8.
For each \(U\) and each \(l\) exactly one of \(\zeta_l\), \(\mathord\sim\zeta_l\)
is satisfied at the given positions, so each \(U\) lies in exactly one
\(C_\varepsilon\): the \(C_\varepsilon\) partition \(Y^{2n+1}\) into \(2^s\) sets,
which is the content of (5.6.4) and (5.6.5)
\cite[pp.~61--62]{EhrenfeuchtMostowski1956}.

By Theorem~5.1 (\cref{thm:p03-w03-p03-028-em-51}) there are \(\varepsilon\) and
an infinite \(Y_1\subseteq Y\) with \(Y_1^{2n+1}\subseteq C_\varepsilon\).
Choose \(y_{-n}\ll\cdots\ll y_n\ll y_{n+1}\) in \(Y_1\).  Both
\(\{y_{-n},\ldots,y_n\}\) and \(\{y_{-n+1},\ldots,y_{n+1}\}\) lie in
\(C_\varepsilon\); the increasing enumeration of the second assigns
\(y_{p+1}\) to position \(p\).  Hence, for each \(l\),
\[
  \operatorname{stsf}_M\zeta_l(y_{i_l},y_{j_l},\ldots,y_{m_l})
  \iff\varepsilon_l=+1\iff
  \operatorname{stsf}_M\zeta_l(y_{i_l+1},y_{j_l+1},\ldots,y_{m_l+1}),
\]
which are (5.6.6) and (5.6.7) \cite[p.~62]{EhrenfeuchtMostowski1956}.  Define
the extension \(M^*\) of \(M\) to \(S^*\) over \(Y\) by \(M^*_{x_j}=y_j\)
(\(-n\leq j\leq n+1\)).  The axioms of \(S\) are valid in \(M^*\) by
\cref{thm:p03-w03-p03-028-em-21}, and the axioms
\(\mathord\sim(x_i\iota x_j)\) because the \(y_j\) are distinct.  Since
\(h(x_j)=x_{j+1}\), Theorem~3.8 (\cref{thm:p03-w03-p03-028-em-38}, applied with
the set \(\{x_{-n},\ldots,x_{n+1}\}\) of new constants) gives
\begin{gather*}
  \operatorname{stsf}_{M^*}\operatorname{subst}\zeta_l(x_{i_l},\ldots,x_{m_l})
  \equiv\operatorname{stsf}_M\zeta_l(y_{i_l},\ldots,y_{m_l}),\\
  \operatorname{stsf}_{M^*}\operatorname{subst}\zeta_l(h(x_{i_l}),\ldots,h(x_{m_l}))
  \equiv\operatorname{stsf}_M\zeta_l(y_{i_l+1},\ldots,y_{m_l+1}),
\end{gather*}
and the two right sides agree.  So the axioms (5.6.1) are valid in \(M^*\),
\(M^*\) is a model of \(S^*\), and \(S^*\) is consistent by the soundness
direction of \cref{thm:p03-w03-p03-028-em-theorems-valid}
\cite[p.~62]{EhrenfeuchtMostowski1956}.
\end{proof}
\end{annexitem}

\begin{annexitem}{annex:p03-w03-p03-028-em-57-fragment}{thm:p03-w03-p03-028-em-57-fragment}{The finite step in the proof of Theorem 5.7}
\begin{proof}
Put \(p=st\).  Let \(X_0=\{x_1,\ldots,x_p\}\), with \(n\) distinct elements
\(x^0_1,\ldots,x^0_n\) (\(n\geq1\)), and \(X_i=\{g_i(x^0_1),\ldots,g_i(x^0_n)\}\)
for \(i\leq s\).  Each \(X_i\) has \(n\) elements, as \(g_i\) is one-one, and
\(X^*=X_0\cup X_1\cup\cdots\cup X_s=\{x^0_1,\ldots,x^0_m\}\) with the
\(x^0_j\) distinct \cite[(5.7.5), p.~64]{EhrenfeuchtMostowski1956}.  Let \(M\)
be a model of \(S\) over an infinite set \(Y\), linearly ordered by \(\ll\)
(axiom of choice).

\emph{Distinguished orderings} \cite[p.~64]{EhrenfeuchtMostowski1956}.  For
\(U\in Y^n\), a sequence \((u_1,\ldots,u_p)\) of elements of \(U\) is a
distinguished ordering of \(U\) if \(u_h\ll u_j\) iff \(x_h\prec x_j\) for
\(h,j\leq p\).  Each \(U\in Y^n\) has exactly one.  Given one, the map
\(x_h\mapsto u_h\) is well defined, since \(x_h=x_j\) excludes both
\(x_h\prec x_j\) and \(x_j\prec x_h\), hence both \(u_h\ll u_j\) and
\(u_j\ll u_h\), so \(u_h=u_j\); it is strictly increasing from \((X_0,\prec)\)
to \((U,\ll)\), hence one-one, hence onto, as both sets have \(n\) elements.
An order isomorphism between two \(n\)-element linear orders sends the
\(k\)-th smallest element to the \(k\)-th smallest, so there is exactly one,
and it defines exactly one distinguished ordering.

For \(\varepsilon\in\{+1,-1\}^s\) let \(C_\varepsilon\) be the set of
\(U\in Y^n\) whose distinguished ordering \((u_1,\ldots,u_p)\) satisfies
\(\operatorname{stsf}_M\zeta_i^{\varepsilon_i}(u_{(i-1)t+1},\ldots,u_{it})\)
for \(i=1,\ldots,s\).  The article states this as (5.7.6) for the matrices
\(\psi_i=\operatorname{subst}\zeta_i(\bar x_{(i-1)t+1},\ldots,\bar x_{it})\),
with distinct variables \(\bar x\) assigned to the elements of \(X^*\)
\cite[p.~64]{EhrenfeuchtMostowski1956}; the formulations agree as in the
previous item.  The \(C_\varepsilon\) partition \(Y^n\) into \(2^s\) disjoint
sets.  By Theorem~5.1 there are \(\varepsilon\) and an infinite
\(Y_1\subseteq Y\) with \(Y_1^n\subseteq C_\varepsilon\)
\cite[pp.~64--65]{EhrenfeuchtMostowski1956}.

Since \(Y_1\) is infinite and \(\prec\) linearly orders \(X^*\), there are
\(y_1,\ldots,y_m\in Y_1\) with \(y_i\ll y_j\) iff \(x^0_i\prec x^0_j\)
\cite[(5.7.7), p.~65]{EhrenfeuchtMostowski1956}: take \(m\) elements of \(Y_1\)
and give \(x^0_i\) the one whose rank among them is the rank of \(x^0_i\) in
\(X^*\).  Let \(M^*\) extend \(M\) to \(S^*\) over \(Y\) by
\(M^*_{x^0_j}=y_j\).  The axioms of \(S\) are valid in \(M^*\) by
\cref{thm:p03-w03-p03-028-em-21}, and the axioms (5.7.2) because the \(y_j\)
are distinct.

Put
\[
  U_0=\{M^*_{x_1},\ldots,M^*_{x_p}\},\qquad
  U_i=\{M^*_{g_i(x_1)},\ldots,M^*_{g_i(x_p)}\}\quad(i\leq s).
\]
For (5.7.8): \(U_0\) has \(n\) elements because \(x\mapsto M^*_x\) is one-one
on \(X^*\), and the sequence of the \(M^*_{x_h}\) (\(h\leq p\)) is its
distinguished ordering, since \(M^*_{x_h}\ll M^*_{x_j}\) iff
\(x_h\prec x_j\) by (5.7.7).  For (5.7.9): \(U_i\) has \(n\) elements, and the
sequence of the \(M^*_{g_i(x_h)}\) is its distinguished ordering; by
(5.7.7) \(M^*_{g_i(x_h)}\ll M^*_{g_i(x_j)}\) iff \(g_i(x_h)\prec g_i(x_j)\), and
since \(g_i\in G(X,\prec)\) this holds iff \(x_h\prec x_j\)
\cite[p.~65]{EhrenfeuchtMostowski1956}.  Both \(U_0\) and \(U_i\) are subsets of
\(Y_1\) with \(n\) elements, so they lie in \(C_\varepsilon\).  Hence for each
\(i\leq s\),
\begin{align*}
  \operatorname{stsf}_M\zeta_i(M^*_{x_{(i-1)t+1}},\ldots,M^*_{x_{it}})
  &\iff\varepsilon_i=+1\\
  &\iff\operatorname{stsf}_M\zeta_i(M^*_{g_i(x_{(i-1)t+1})},\ldots,M^*_{g_i(x_{it})}).
\end{align*}
By Theorem~3.8 (with the set \(X^*\) of new constants) the two sides are the
truth values in \(M^*\) of the two closed matrices
\[
  \operatorname{subst}\zeta_i(x_{(i-1)t+1},\ldots,x_{it}),\qquad
  \operatorname{subst}\zeta_i(g_i(x_{(i-1)t+1}),\ldots,g_i(x_{it})),
\]
so the axioms (5.7.3) are valid in \(M^*\) \cite[(5.7.10), p.~66]{EhrenfeuchtMostowski1956}.
The article reaches the second side by passing from \(\psi_i\) to the matrix
\(\bar\psi_i\), obtained by renaming the variables \(\bar x_k\) to
\(\overline{g_i(x_k)}\), and inserting variables that do not occur
\cite[pp.~64, 66]{EhrenfeuchtMostowski1956}; the statement above is what that
passage establishes.  So \(M^*\) is a model of \(S^*\), and \(S^*\) is
consistent.
\end{proof}
\end{annexitem}

\begin{annexitem}{annex:p03-w03-p03-028-em-open-reduction}{thm:p03-w03-p03-028-em-open-reduction}{Reduction to open theories}
\begin{proof}
Let \(M\) be a model of \(S\) over an infinite set.  For \(k\geq1\) let
\(\delta_k\) be the closed matrix of \(S\), built from \(\iota\) alone, saying
that there are at least \(k\) elements.  Each \(\delta_k\) is valid in \(M\), so
\(\mathord\sim\delta_k\notin T(S)\) by
\cref{thm:p03-w03-p03-028-em-theorems-valid}; as
\(\mathord\sim\delta_k\in Z(S)\) and \(T(S')\cap Z(S)=T(S)\),
\(\mathord\sim\delta_k\notin T(S')\).  Suppose the closures of the axioms of
\(S'\) together with all \(\delta_k\) derived \(\bottom\).  The derivation uses
finitely many \(\delta_k\), all derivable from \(\delta_K\) for the largest
\(K\) among them, so the closures of the axioms of \(S'\) and \(\delta_K\)
derive \(\bottom\), and by \(\neg\)-introduction \(\mathord\sim\delta_K\in
T(S')\), a contradiction.  So this set is consistent, and by completeness
(\cref{thm:fol-completeness}) it has a model, which satisfies every
\(\delta_k\) and is therefore infinite; it is a model of \(S'\) over an
infinite set.

Let \(M'\) be a model of \(S'\) over \(X_1\) and \(M\) the pseudo-model of
\(S\) over \(X_1\) that agrees with \(M'\) on \(F(S)\cup P(S)\).  By
\cref{thm:p03-w03-p03-028-em-theorems-valid}, \(A(S)\subseteq T(S)\subseteq
T(S')\subseteq V_{M'}\), and by \cref{thm:p03-w03-p03-028-em-21}
\(V_{M'}\cap Z(S)=V_M\); so \(M\) is a model of \(S\).  By Lemma~4.2
(\cref{thm:p03-w03-p03-028-em-41}), \(G_{M'}\subseteq G_M\).  If every element
of a group \(G\) of transformations of \(X\subseteq X_1\) extends to an
element of \(G_{M'}\), that element lies in \(G_M\), so \(G_M\) strongly
contains \(G\).

For Theorem~5.6, apply its open case to \(S'\), which has a model over an
infinite set: some model \(M'_0\) of \(S'\) has \(G_{M'_0}\) strongly
containing an infinite cyclic group, and so does the corresponding model of
\(S\).  For Theorem~5.7 the same argument applies for each ordered set
\((X,\prec)\) \cite[pp.~62--63]{EhrenfeuchtMostowski1956}.
\end{proof}
\end{annexitem}

\begin{annexitem}{annex:p03-w03-p03-028-em-61}{thm:p03-w03-p03-028-em-61}{Ehrenfeucht--Mostowski Theorem 6.1}
\begin{proof}
If \(X\) is empty, take \(M_0=M\) and \(X_1=Y\): (6.1.1) holds because the only
transformation of \(\varnothing\) is extended by the identity, and (6.1.2) is
empty when \(q\geq1\), while for \(q=0\) the closed matrix \(\eta\) is
satisfied in \(M\) because the empty sequence is increasing in \(Y^*\).  Let
\(X\) be non-empty.

Let \(S_1\) be \(S^*(X)\) with these axioms added
\cite[p.~67]{EhrenfeuchtMostowski1956}: all equivalences
(5.4.1) with \(h\in G(X,\prec)\), \(\zeta\) open and
\(x_1,x_2,\ldots\in X\); and all matrices
\(\operatorname{subst}\eta(x_1,\ldots,x_q)\) with
\(x_1\prec\cdots\prec x_q\) in \(X\) (6.1.3), the variables \(\xi_j\) with
\(j>q\) not occurring in \(\eta\).

\emph{\(S_1\) is consistent.}  A derivation of \(\bottom\) would use the closures
of finitely many axioms \(\Phi\) of \(S_1\).  List the equivalences of type
(5.4.1) in \(\Phi\) as \(i=1,\ldots,s\), each with its
open matrix \(\zeta_i\), its \(g_i\in G(X,\prec)\) and its assignment; let
\(t\geq1\) bound the indices of the variables in all \(\zeta_i\), and let
\(x_{(i-1)t+1},\ldots,x_{it}\) be the values the \(i\)-th assignment gives to
\(\xi_1,\ldots,\xi_t\).  For each element \(x\) of \(X\) that occurs in
\(\Phi\) but not yet among these values, add an instance with
\(\zeta=\xi_1\iota\xi_1\), \(g\) the identity of \(X\) and every value equal to
\(x\); if \(\Phi\) has no equivalences and no elements of \(X\), add one such
instance for some \(x\in X\).  Let \(X^*\) and \(S^*\) be as in
\cref{thm:p03-w03-p03-028-em-57-fragment} for these data, and add to \(S^*\) the
matrices (6.1.3) with \(x_1,\ldots,x_q\in X^*\).  Every member of \(\Phi\) is an
axiom of this theory, which the article calls \(S^*\)
\cite[p.~67]{EhrenfeuchtMostowski1956}.  Repeat the proof of
\cref{thm:p03-w03-p03-028-em-57-fragment} with one change: the partition is of
\(Y^{*n}\) instead of \(Y^n\), and Theorem~5.1 is applied to the infinite set
\(Y^*\), so that \(Y_1\subseteq Y^*\) and \(M^*_x\in Y^*\) for \(x\in X^*\)
\cite[p.~67]{EhrenfeuchtMostowski1956}.  The axioms of \(S\), (5.7.2) and
(5.7.3) are valid in \(M^*\) as before.  If \(x_1,\ldots,x_q\in X^*\) and
\(x_1\prec\cdots\prec x_q\), then \(M^*_{x_1}\ll\cdots\ll M^*_{x_q}\) by
(5.7.7), all in \(Y^*\), so
\(\operatorname{stsf}_M\eta(M^*_{x_1},\ldots,M^*_{x_q})\) by hypothesis, and by
Theorem~3.8 \(\operatorname{subst}\eta(x_1,\ldots,x_q)\) is valid in \(M^*\).
So the theory has a model and \(\Phi\) is consistent, a contradiction.

\emph{The model.}  By \cref{thm:p03-w03-p03-028-em-complete-set} there is a
consistent and complete \(I\supseteq A(S_1)\); it contains \(A(S^*(X))\), the
equivalences (5.4.1) and the matrices (6.1.3)
\cite[pp.~67--68]{EhrenfeuchtMostowski1956}.  \(\mathfrak M(X,I)\) is a model of
\(S\) (\cref{thm:p03-w03-p03-028-em-37}); each \(h\in G(X,\prec)\) induces
\(f_h\) (\cref{thm:p03-w03-p03-028-em-54}); and for
\(x_1\prec\cdots\prec x_q\), Lemma~3.6 with \(\tau_j=x_j\) (\(j\leq q\)) and
\(\tau_j=x_1\) (\(j>q\)) gives
\(\operatorname{stsf}_{\mathfrak M(X,I)}\eta([x_1],\ldots,[x_q])\), because
\(\operatorname{subst}\eta(x_1,\ldots,x_q)\in I\).  Rename as in the proof of
\cref{thm:p03-w03-p03-028-em-55}, identifying \([x]\) with \(x\)
\cite[p.~68]{EhrenfeuchtMostowski1956}: the renamed model \(M_0\) over
\(X_1\supseteq X\) satisfies \(\eta\) at every increasing \(q\)-tuple of \(X\),
since renaming preserves satisfaction, and
\(\beta f_h\beta^{-1}\in G_{M_0}\) extends \(h\).  This gives (6.1.1) and
(6.1.2).
\end{proof}
\end{annexitem}

The last three items prove the results of
\cref{sec:p03-w03-p03-028-two-models}.

\begin{annexitem}{annex:p03-w03-p03-028-thin-tuples}{thm:p03-w03-p03-028-thin-tuples}{Ehrenfeucht--Mostowski models over well-orders realize few types}
\begin{proof}
\emph{The data exist.}  Let \(\mathcal A\) be an infinite model of \(T\).
Interpreting each \(R_\varphi\) by the set defined by \(\varphi\) makes
\(\mathcal A\) a model of \(T^*\), and by
\cref{def:p03-morley-skolem-hull} this model has an expansion to a model of
\(T^{*\mathrm{sk}}\), which is still infinite.  By
\cref{thm:p03-morley-ramsey-em}, applied to \(T^{*\mathrm{sk}}\) and the
ordered set \(\omega\), some model of \(T^{*\mathrm{sk}}\) contains
distinct elements \((a_i)_{i<\omega}\) that are order indiscernible, in
particular for quantifier-free formulas, so \(\Phi\) of
\cref{thm:p03-w03-p03-028-coherent} exists.  By
\cref{thm:p03-w03-p03-028-realize} with \(J=\kappa\) there is a sequence
\((b_j)_{j<\kappa}\) realizing \(\Phi\) in a model \(\mathcal K\) of
\(T^{*\mathrm{sk}}\).

\emph{\(\mathcal M\) is a model of \(T\) of cardinality \(\kappa\).}  By
\cref{thm:p03-w03-p03-028-realize}, applied to \(T^*\), \(H_\kappa\) is a
model of \(T^{*\mathrm{sk}}\); its reduct to the language of \(T^*\) is a
model of \(T^*\), whose axioms include those of \(T\), so
\(\mathcal M\entails T\).  The language \(\Lang\) is countable, hence so
are the Morleyized language and its Skolem expansion
(\cref{def:p03-morley-skolem-hull}), and there are countably many terms.
Each element of \(H_\kappa\) is the value of a term at a finite increasing
tuple from \(\kappa\); there are \(\kappa\) such tuples, so
\(\abs{\dom\mathcal M}\leq\aleph_0\cdot\kappa=\kappa\).  The \(b_j\) are
distinct, because \(p_2\) contains \(v_1\neq v_2\), so
\(\abs{\dom\mathcal M}\geq\kappa\).

\emph{Representations.}  Let \(Y\subseteq\dom\mathcal M\) be countable and
\(n\geq1\).  For each \(y\in Y\) fix a term \(s_y\) and a finite increasing
tuple \(\bar x_y\) from \(\kappa\) with \(y=s_y(b_{\bar x_y})\), and let
\(X_0\subseteq\kappa\) be the union of the sets of entries of the
\(\bar x_y\).  As a countable union of finite sets, \(X_0\) is countable.
For \(j<\kappa\) let \(\ell(j)\) be the least element of
\(\{x\in X_0:x\geq j\}\) if this set is nonempty, and \(\ell(j)=\infty\)
otherwise; the least element exists because \(\kappa\) is well-ordered.
Then, for \(x\in X_0\),
\begin{equation*}
 j\leq x\iff \ell(j)\leq x,
 \qquad
 j=x\iff \ell(j)=x\text{ and }j\in X_0,
\end{equation*}
where \(\infty\leq x\) is false: if \(j\leq x\) then \(x\) is in the set and
\(\ell(j)\leq x\); if \(\ell(j)\leq x\) then \(j\leq\ell(j)\leq x\).  The
second equivalence holds because \(j\in X_0\) gives \(\ell(j)=j\).

For an \(n\)-tuple \(\bar c\) of \(\mathcal M\) fix terms
\(t_1,\ldots,t_n\) in variables \(v_1,\ldots,v_r\) and indices
\(j_1<\cdots<j_r\) in \(\kappa\) with
\(c_k=t_k(b_{j_1},\ldots,b_{j_r})\) for \(k=1,\ldots,n\); a common
increasing tuple exists because the union of finitely many finite sets of
indices can be listed increasingly and each term can be read as a term in
all the listed variables.  The \emph{pattern} of \(\bar c\) is the tuple
consisting of \(r\), the terms \(t_1,\ldots,t_n\), and for each
\(s\leq r\) the pair \((\ell(j_s),\epsilon_s)\), where \(\epsilon_s=1\) if
\(j_s\in X_0\) and \(\epsilon_s=0\) otherwise.

\emph{Countably many patterns.}  There are countably many choices of
\(r\), countably many of the terms, and for each \(s\) countably many pairs
from \((X_0\cup\{\infty\})\times\{0,1\}\).  So there are countably many
patterns.

\emph{Equal patterns give equal types.}  Let \(\bar c\) and \(\bar c'\)
have the same pattern, with indices \(j_1<\cdots<j_r\) and
\(j'_1<\cdots<j'_r\) and the same terms.  Let
\(I=X_0\cup\{j_1,\ldots,j_r\}\) with the order of \(\kappa\), and define
\(e\colon I\to\kappa\) by \(e(x)=x\) for \(x\in X_0\) and \(e(j_s)=j'_s\).
This is well defined: if \(j_s\in X_0\), then \(\epsilon_s=1\), so
\(j'_s\in X_0\) and \(j'_s=\ell(j'_s)=\ell(j_s)=j_s\).  It preserves the
order strictly: on \(X_0\) it is the identity; \(j_s<j_{s'}\) iff
\(s<s'\) iff \(j'_s<j'_{s'}\); and for \(x\in X_0\) the relations
\(j_s\leq x\) and \(j_s=x\) are determined by \((\ell(j_s),\epsilon_s)\)
through the displayed equivalences, which are the same for \(j'_s\).  So
\(e\) is an order embedding, and so is the inclusion
\(\iota\colon I\to\kappa\).

The subsequence \((b_i)_{i\in I}\) realizes \(\Phi\), since its increasing
tuples are increasing tuples of \((b_j)_{j<\kappa}\), and its term hull
\(H_I\) is a model of \(T^{*\mathrm{sk}}\) by the argument of
\cref{thm:p03-w03-p03-028-realize}.  Apply
\cref{thm:p03-w03-p03-028-elementary} to the order embeddings \(\iota\)
and \(e\) of \(I\) into \(\kappa\), with \(H_J=H_\kappa\).  The map
\(\hat\iota\) induced by \(\iota\) is the inclusion of \(H_I\) in
\(H_\kappa\), and the map \(\hat e\) induced by \(e\) sends
\(t(b_{i_1},\ldots,b_{i_p})\) to \(t(b_{e(i_1)},\ldots,b_{e(i_p)})\); both
are elementary on the \(\Lang\)-reducts.  Every \(y\in Y\) lies in \(H_I\),
and \(\hat e(y)=s_y(b_{e(\bar x_y)})=s_y(b_{\bar x_y})=y\) because \(e\)
fixes \(X_0\).  The tuple \(\bar c\) lies in \(H_I\), and
\(\hat e(c_k)=t_k(b_{j'_1},\ldots,b_{j'_r})=c'_k\).  Hence for every
\(\Lang\)-formula \(\psi(\bar x,\bar y)\) and every tuple \(\bar y\) from
\(Y\),
\[
 \mathcal M\entails\psi(\bar c,\bar y)
 \iff H_I\entails\psi(\bar c,\bar y)
 \iff \mathcal M\entails\psi(\hat e\bar c,\hat e\bar y)
 \iff \mathcal M\entails\psi(\bar c',\bar y),
\]
the first equivalence by elementarity of \(\hat\iota\), the second by
elementarity of \(\hat e\).  So
\(\operatorname{tp}(\bar c/Y)=\operatorname{tp}(\bar c'/Y)\).

Every realized \(n\)-type over \(Y\) is the type of a tuple with some
pattern, and the type depends only on the pattern, so at most countably
many \(n\)-types over \(Y\) are realized in \(\mathcal M\).
\end{proof}
\end{annexitem}

\begin{annexitem}{annex:p03-w03-p03-028-thick}{thm:p03-w03-p03-028-thick}{A model with many types over a countable set}
\begin{proof}
By \cref{def:p03-stability} there are \(\mathcal A\entails T\) and tuples
\((a_i)_{i<\omega}\), \((b_i)_{i<\omega}\) of \(\mathcal A\) with
\(\mathcal A\entails\varphi(a_i;b_l)\) iff \(i<l\).  The \(a_i\) are
pairwise distinct: for \(i<i'\), \(\varphi(a_i;b_{i'})\) holds and
\(\varphi(a_{i'};b_{i'})\) fails.  So \(\mathcal A\) is infinite.  The set
\(\mathbb R\setminus\mathbb Q\) has cardinality \(2^{\aleph_0}\geq\aleph_1\),
so it has a subset \(R\) of cardinality \(\aleph_1\).  Let
\(D=\mathbb Q\cup R\), ordered as a subset of \(\mathbb R\).

Add new constants: tuples \(\bar c_d\) of length \(n\) and \(\bar e_d\) of
the length of \(\bar y\) for \(d\in D\), and constants \(g_\alpha\) for
\(\alpha<\kappa\).  Let \(\Theta\) consist of \(T\), the sentences
\(\varphi(\bar c_d;\bar e_{d'})\) for \(d<d'\) in \(D\), the sentences
\(\neg\varphi(\bar c_d;\bar e_{d'})\) for \(d\geq d'\) in \(D\), and
\(g_\alpha\neq g_\beta\) for \(\alpha<\beta<\kappa\).  A finite
\(\Theta_0\subseteq\Theta\) mentions finitely many \(d\), say
\(d_1<\cdots<d_p\), and finitely many \(g_\alpha\).  In \(\mathcal A\)
interpret \(\bar c_{d_k}\) as \(a_k\), \(\bar e_{d_k}\) as \(b_k\), and the
mentioned \(g_\alpha\) as distinct elements, which exist because
\(\mathcal A\) is infinite.  Then \(\varphi(\bar c_{d_k};\bar e_{d_m})\) is
interpreted as \(\varphi(a_k;b_m)\), which holds iff \(k<m\), iff
\(d_k<d_m\).  So \(\Theta_0\) has a model, and by compactness
(\cref{thm:fol-compactness}) \(\Theta\) has a model.  Let \(\mathcal N'\)
be its \(\Lang\)-reduct, and keep writing \(\bar c_d\), \(\bar e_d\) and
\(g_\alpha\) for the interpretations of the new constants.  Then
\(\mathcal N'\entails T\) and \(\abs{\dom\mathcal N'}\geq\kappa\).

Let \(A_0\) be the set of entries of all \(\bar c_d\) and \(\bar e_d\),
which has cardinality at most \(\aleph_1\leq\kappa\), and let \(B\) be the
set of the \(g_\alpha\), of cardinality \(\kappa\).  By
\cref{thm:downward-ls} there is \(\mathcal N\preccurlyeq\mathcal N'\)
containing \(A_0\cup B\) with
\(\abs{\dom\mathcal N}\leq\max(\kappa,\aleph_0)=\kappa\); since it contains
\(B\), \(\abs{\dom\mathcal N}=\kappa\), and \(\mathcal N\entails T\).

Let \(Y\) be the set of entries of the tuples \(\bar e_q\) for
\(q\in\mathbb Q\); it is countable.  Let \(r<r'\) be in \(R\).  There is a
rational \(q\) with \(r<q<r'\).  By the choice of \(\Theta\),
\(\mathcal N'\entails\varphi(\bar c_r;\bar e_q)\) and
\(\mathcal N'\entails\neg\varphi(\bar c_{r'};\bar e_q)\), and the same holds
in \(\mathcal N\) because \(\mathcal N\preccurlyeq\mathcal N'\) and the
tuples lie in \(\mathcal N\).  So the formula \(\varphi(\bar x;\bar e_q)\),
with parameters from \(Y\), belongs to
\(\operatorname{tp}(\bar c_r/Y)\) and not to
\(\operatorname{tp}(\bar c_{r'}/Y)\).  The \(n\)-types
\(\operatorname{tp}(\bar c_r/Y)\), \(r\in R\), are therefore pairwise
distinct, and there are \(\aleph_1\) of them.
\end{proof}
\end{annexitem}

\begin{annexitem}{annex:p03-w03-p03-028-two-models}{thm:p03-w03-p03-028-two-models}{Unstable countable theories are not categorical in uncountable cardinals}
\begin{proof}
Let \(\varphi(\bar x;\bar y)\) have the order property in \(T\), with
\(\bar x\) of length \(n\).  As shown in the proof of
\cref{thm:p03-w03-p03-028-thick}, \(T\) has an infinite model, so
\cref{thm:p03-w03-p03-028-thin-tuples} applies.  Both \(\mathcal M\) and
\(\mathcal N\) are models of \(T\) of cardinality \(\kappa\).

Suppose \(F\colon\mathcal N\to\mathcal M\) is an isomorphism.  Let \(Y\)
and the tuples \(\bar c_r\) (\(r\in R\)) be as in the proof of
\cref{thm:p03-w03-p03-028-thick}, with pairwise distinct types over
\(Y\).  For \(r\neq r'\) there are a formula \(\psi(\bar x,\bar y)\) and a
tuple \(\bar y\) from \(Y\) with \(\mathcal N\entails\psi(\bar c_r,\bar
y)\) and \(\mathcal N\entails\neg\psi(\bar c_{r'},\bar y)\): two distinct
types differ in some formula, and each type contains every formula with
parameters from \(Y\) or its negation.  By \cref{lem:iso-invariance},
\(\mathcal M\entails\psi(F\bar c_r,F\bar y)\) and
\(\mathcal M\entails\neg\psi(F\bar c_{r'},F\bar y)\).  So the tuples
\(F\bar c_r\), \(r\in R\), realize \(\aleph_1\) pairwise distinct
\(n\)-types over the countable set \(F[Y]\subseteq\dom\mathcal M\), which
contradicts \cref{thm:p03-w03-p03-028-thin-tuples}.  Hence \(\mathcal M\)
and \(\mathcal N\) are not isomorphic, and \(T\) is not
\(\kappa\)-categorical.
\end{proof}
\end{annexitem}
